\documentclass[14pt,a4paper,oneside,openany]{extbook}
\let\cleardoublepage\clearpage

% XeLaTeX: шрифты через fontspec, русский язык через polyglossia.
% Конфигурация по канону Russian-Phd-LaTeX-Dissertation-Template (А. Акиньшин):
% явно объявляется \cyrillicfont для polyglossia, иначе переносы и метрики кириллицы
% подхватываются нестабильно.
\usepackage{fontspec}
\defaultfontfeatures{Ligatures=TeX}

\usepackage{polyglossia}
\setmainlanguage{russian}
\setotherlanguage{english}

\setmainfont{Times New Roman}
\newfontfamily\cyrillicfont{Times New Roman}
\newfontfamily\englishfont{Times New Roman}

% Моноширинный шрифт с поддержкой кириллицы (для \texttt{...} рядом с кириллическим текстом).
\setmonofont{Courier New}
\newfontfamily\cyrillicfonttt{Courier New}

% Математические пакеты (\mathbb, \varnothing, \mathcal, окружения align/equation и т.п.)
\usepackage{amsmath}
\usepackage{amssymb}

% Русская типография: одиночный пробел после точки (вместо English-style увеличенного).
\frenchspacing

% Умеренная защита от overfull \hbox. Жёсткие настройки не нужны:
% переносы корректно работают через polyglossia с русскими паттернами.
\emergencystretch=1em

% Уменьшение вертикальных отступов до и после display-формул:
% по умолчанию у extbook 14pt они слишком велики, а \onehalfspacing их дополнительно растягивает.
\setlength{\abovedisplayskip}{8pt plus 2pt minus 2pt}
\setlength{\belowdisplayskip}{8pt plus 2pt minus 2pt}
\setlength{\abovedisplayshortskip}{4pt plus 2pt minus 2pt}
\setlength{\belowdisplayshortskip}{4pt plus 2pt minus 2pt}

% Настройка полей: левое 30 мм, правое 15 мм (допускается 10 мм), верхнее 20 мм, нижнее 20 мм (п. 4.2)
\usepackage[a4paper,top=20mm,bottom=20mm,left=30mm,right=15mm]{geometry}

% Межстрочный интервал 1.5, размер шрифта 14 пт (п. 4.2)
\usepackage{setspace}
\onehalfspacing

% Абзацный отступ 1.25 см (п. 4.2)
\setlength{\parindent}{1.25cm}
\setlength{\parskip}{0pt}

% Нумерация страниц: арабскими цифрами, по центру нижней части, без точки (п. 4.3.1)
% Нумерация сквозная, включая приложения — см. \setcounter{page} в main.tex
\usepackage{fancyhdr}
\pagestyle{fancy}
\fancyhf{}
\renewcommand{\headrulewidth}{0pt}
\cfoot{\thepage}
\fancypagestyle{plain}{%
  \fancyhf{}%
  \renewcommand{\headrulewidth}{0pt}%
  \cfoot{\thepage}%
}

% Настройка заголовков через titlesec
\usepackage{titlesec}
% Главы: прописными, полужирный, по центру, без точки (п. 4.4.1, 4.4.4)
\titleformat{\chapter}[hang]{\bfseries\fontsize{16pt}{18pt}\selectfont\centering}{\thechapter}{1em}{\MakeUppercase}
% Разделы: полужирный, с абзацного отступа (п. 4.4.4)
\titleformat{\section}[hang]{\bfseries\normalsize}{\thesection}{1em}{}
% Подразделы: обычный шрифт, с абзацного отступа (п. 4.4.4 — полужирный только для разделов)
\titleformat{\subsection}[hang]{\normalfont}{\thesubsection}{1em}{}
\titlespacing{\chapter}{0pt}{0pt}{20pt}
\titlespacing{\section}{1.25cm}{12pt}{6pt}
\titlespacing{\subsection}{1.25cm}{6pt}{3pt}

% Глубина содержания и нумерации
\setcounter{tocdepth}{2}      % в содержании: главы + разделы + подразделы
\setcounter{secnumdepth}{3}   % нумерация до уровня подпунктов (1.1.1)

% Настройка содержания: обычный шрифт (не жирный) (п. 4.2), точечные лидеры (п. 3.4)
\usepackage{tocloft}
\renewcommand{\cftdotsep}{1}
\renewcommand{\cftchapleader}{\cftdotfill{\cftdotsep}}
\renewcommand{\cftchapfont}{\normalfont}
\renewcommand{\cftchappagefont}{\normalfont}
\renewcommand{\cftsecfont}{\normalfont}
\renewcommand{\cftsecpagefont}{\normalfont}
\renewcommand{\cftsubsecfont}{\normalfont}
\renewcommand{\cftsubsecpagefont}{\normalfont}

% Конфигурационные файлы
\usepackage{enumitem}
\setlist{nosep} % Убирает дополнительные отступы между элементами списков

% По ГОСТ 7.32 и требованиям ИТМО (Приложение 1) маркированные списки — через тире.
% Маркер выставляется с абзацного отступа 1.25 см (норма ИТМО п. 4.5).
\setlist[itemize]{leftmargin=1.25cm, label=--}

% Нумерованные списки: арабские цифры со скобкой: 1), 2), 3) ...
\setlist[enumerate]{leftmargin=1.25cm, label=\arabic*)}

\input{config-formulas.tex}
% Подключаем нужные пакеты для таблиц и оформления
\usepackage{booktabs}   % Для красивых горизонтальных линий (\toprule, \midrule, \bottomrule)
\usepackage{multirow}   % Для объединения ячеек по вертикали
\usepackage{caption}    % Для \ContinuedFloat и настройки заголовков
\usepackage{longtable}
\usepackage{etoolbox}

\DeclareCaptionLabelSeparator{emdash}{\space---\space}  % пробел, длинное тире, пробел

% Настройка заголовков таблиц: слева без абзацного отступа, разделитель — тире (п. 4.6.3)
\captionsetup[table]{position=t, singlelinecheck=false, justification=raggedright, labelsep=emdash, name=Таблица}

% Переключатель для смены выравнивания заголовков при переносе таблицы (п. 4.6.3)
\newcommand*\RightCaption{\captionsetup{justification=raggedleft}}
\newcommand*\LeftCaption{\captionsetup{justification=raggedright}}

% ------------------------
% Шаблон: Многоуровневый заголовок
% ------------------------
\newcommand{\ExampleMultiLevelTable}{%
\begin{table}[h!]
\caption{Пример построения таблицы}
\centering
\begin{tabular}{|c|c|c|c|}
\hline
\multirow{2}{*}{\textbf{Заголовок}} & \multicolumn{2}{c|}{\textbf{Заголовок колонки}} & \multirow{2}{*}{\textbf{Заголовок колонки}}\\
\cline{2-3}
& \textbf{Подзаголовок} & \textbf{Подзаголовок} & \\
\hline
A & B & C & D \\
\hline
\end{tabular}
\end{table}
}

% Шаблон: «Разбитая» таблица (две части под одним номером)
% Номер таблицы в строке «Продолжение таблицы N» подставляется автоматически
\newcommand{\ExampleSplitTable}{%
  % --- Первая часть ---
  \begin{table}[h!]
    \LeftCaption
    \caption{Пример переноса таблицы}
    \label{tab:split_example}
    \centering
    \begin{tabular}{|c|c|c|c|}
    \hline
    \textbf{Заголовок} & \textbf{Заголовок колонки} & \textbf{Заголовок колонки} & \textbf{Заголовок колонки} \\
    \hline
    1 & 2 & 3 & 4 \\
    \hline
    2 & 3 & 4 & 5 \\
    \hline
    \end{tabular}
  \end{table}

  % --- Продолжение (номер подставляется автоматически через \thetable) ---
  \begin{table}[h!]
    \ContinuedFloat
    \RightCaption
    \caption*{Продолжение таблицы \thetable}
    \centering
    \begin{tabular}{|c|c|c|c|}
    \hline
    \textbf{Заголовок} & \textbf{Заголовок колонки} & \textbf{Заголовок колонки} & \textbf{Заголовок колонки} \\
    \hline
    3 & 4 & 5 & 6 \\
    \hline
    4 & 5 & 6 & 7 \\
    \hline
    \end{tabular}
  \end{table}
}

\usepackage{graphicx}
\graphicspath{{images/}} % Папка с изображениями
% caption и emdash-разделитель уже объявлены в config-tables.tex

% Настройка подписей к рисункам: по центру, под рисунком, разделитель — тире (п. 4.5.3)
\captionsetup[figure]{position=b, singlelinecheck=false, labelsep=emdash, justification=centering, name=Рисунок}


% Библиография: ГОСТ-numeric, порядок по первому цитированию (п. 4.10.3–4.10.5)
\usepackage[backend=biber,style=gost-numeric,sorting=none]{biblatex}
\addbibresource{bibliography.bib}

% Убираем курсив из названий источников
\DeclareFieldFormat*{title}{#1}
\DeclareFieldFormat*{journaltitle}{#1}
\DeclareFieldFormat*{booktitle}{#1}
\DeclareFieldFormat*{issuetitle}{#1}
\DeclareFieldFormat{namewrapperformat}{#1}

% Убираем курсив из имён авторов (\mkgostheading — внутренняя команда gost-numeric)
\renewcommand*{\mkgostheading}[1]{#1}
\renewcommand*{\mkbibnamefamily}[1]{#1}
\renewcommand*{\mkbibnamegiven}[1]{#1}
\renewcommand*{\mkbibnameprefix}[1]{#1}
\renewcommand*{\mkbibnamesuffix}[1]{#1}

% Буквы для приложений (исключены Ё, З, Й, О, Ч, Ъ, Ы, Ь — п. 4.11.2)
\def\cyrtoc#1{\ifcase #1\or А\or Б\or В\or Г\or Д\or Е\or Ж\or И\or К\or Л\or М\or Н\or П\or Р\or С\or Т\or У\or Ф\or Х\or Ц\or Ш\or Щ\or Э\or Ю\or Я\else \@ctrerr \fi}

\newcommand\originalappendix{}
\let\originalappendix\appendix
\renewcommand{\appendix}{%
  \originalappendix
  \renewcommand{\thechapter}{\cyrtoc{\value{chapter}}}
  \renewcommand{\appendixname}{ПРИЛОЖЕНИЕ}
  % Заголовок приложения: «ПРИЛОЖЕНИЕ А» по центру, затем название (п. 4.11.2)
  \titleformat{\chapter}{\fontsize{16pt}{18pt}\selectfont\bfseries\centering}{\MakeUppercase{\chaptertitlename}.\ \thechapter}{4pt}{}
  % В приложениях оглавление сокращается до уровня глав
  \addtocontents{toc}{\protect\value{tocdepth}=0}
}

% Переопределяем \mainmatter: убираем сброс счётчика страниц (п. 4.3.1 — сквозная нумерация)
\makeatletter
\renewcommand{\mainmatter}{%
  \cleardoublepage
  \@mainmattertrue
}
\makeatother


\begin{document}

% Нумерация страниц начинается с листа «СОДЕРЖАНИЕ» (п. 4.3.1)
% Нумерация сквозная по всему документу, включая приложения.
% Титульный лист, задание и аннотация оформляются в ИСУ ИТМО отдельно.
% Если нужно скорректировать стартовый номер (например, 4 если перед
% содержанием идут 3 страницы из ИСУ) — измените число ниже:
\setcounter{page}{1}

% СОДЕРЖАНИЕ — отдельный лист, без номера раздела (п. 4.4.1, 3.4)
\chapter*{СОДЕРЖАНИЕ}
\markboth{СОДЕРЖАНИЕ}{}
\vspace{-3.5cm}
{%
  \renewcommand{\contentsname}{}%
  \tableofcontents
}

% Список сокращений и условных обозначений (п. 3.2, 4.8)
\chapter*{СПИСОК СОКРАЩЕНИЙ И УСЛОВНЫХ ОБОЗНАЧЕНИЙ}
\addcontentsline{toc}{chapter}{СПИСОК СОКРАЩЕНИЙ И УСЛОВНЫХ ОБОЗНАЧЕНИЙ}
\begingroup
\setlength{\parindent}{0pt}
% Перечень располагается столбцом, слева без абзацного отступа в алфавитном порядке,
% справа через тире — расшифровка, без знаков препинания в конце строки (п. 4.8.4)
ВКР --- выпускная квалификационная работа\\[0.5em]
СППР --- система поддержки принятия решений\\[0.5em]
LHS --- латинский гиперкуб (Latin Hypercube Sampling)\\[0.5em]
LLM --- большая языковая модель (Large Language Model)\\[0.5em]
RDM --- метод робастного принятия решений (Robust Decision Making)
\endgroup

\cleardoublepage

% Список терминов и определений (п. 3.2, 4.9)
\chapter*{ТЕРМИНЫ И ОПРЕДЕЛЕНИЯ}
\addcontentsline{toc}{chapter}{ТЕРМИНЫ И ОПРЕДЕЛЕНИЯ}
\begingroup
\setlength{\parindent}{0pt}
% Перечень располагается столбцом, слева без абзацного отступа в алфавитном порядке,
% справа через тире — определение, без знаков препинания в конце строки (п. 4.9.2)
% Термины пишутся без жирного шрифта
Граничная верификация модели --- проверка имитационной модели на удовлетворение обязательным свойствам в граничных значениях её параметров и в режимах редукции\\[0.5em]
Косинусная близость --- мера сходства двух векторов, определяемая как косинус угла между ними\\[0.5em]
Латинский гиперкуб --- план эксперимента, обеспечивающий равномерное стратифицированное покрытие пространства параметров заданным числом точек при сокращении объёма по сравнению с полным факторным перебором\\[0.5em]
Общие случайные числа --- приём планирования имитационных экспериментов, при котором сравниваемые альтернативы выполняются на одной и той же реализации источников случайности для снижения шума попарного сравнения\\[0.5em]
Политика рекомендаций --- модуль, отображающий программу конференции и профиль участника в упорядоченный список рекомендуемых докладов в текущем временном слоте\\[0.5em]
Социальное заражение --- компонент модели поведения участника, отражающий монотонное усиление склонности к выбору доклада с ростом числа уже сделанных в когорте выборов того же доклада\\[0.5em]
Сценарный анализ --- подход к принятию решений в условиях глубокой неопределённости, в котором стратегии оцениваются на множестве правдоподобных сценариев, а не на единственной точечной прогнозной модели\\[0.5em]
Функция полезности --- числовая функция, по значениям которой формируется распределение вероятностей выбора участника между доступными в слоте докладами
\endgroup

\cleardoublepage

% Основной текст работы (п. 3.2)
\mainmatter
\chapter*{ВВЕДЕНИЕ}
\addcontentsline{toc}{chapter}{ВВЕДЕНИЕ}
\sloppy


\section*{Актуальность темы}
\addcontentsline{toc}{section}{Актуальность темы}

Современные научно-практические IT-конференции — многодневные мероприятия с десятками параллельных сессий, в которых участвуют тысячи специалистов. Программа конференции включает доклады разной тематики, уровня и популярности; залы существенно различаются по вместимости. Распределение участников между параллельными сессиями определяет содержательную ценность мероприятия для участников, репутацию и коммерческую привлекательность мероприятия для организаторов и охват аудитории для докладчиков.

В отсутствие координирующего вмешательства распределение формируется как агрегат независимых индивидуальных выборов. Это приводит к одновременной перегрузке популярных залов, которые не могут вместить желающих, и хронической недозагрузке менее популярных, что снижает суммарное качество мероприятия. Прямое применение классических рекомендательных систем не решает задачи: оптимизируя индивидуальную релевантность независимо для каждого пользователя, такие системы воспроизводят одинаково популярные позиции в верхних строках выдачи и тем самым усиливают, а не ослабляют нагрузку на наиболее востребованные ресурсы.

Описанная задача относится к более широкому классу задач распределения автономных пользователей по ресурсам ограниченной вместимости через ненавязчивую рекомендательную подсказку. Этот класс задач возникает в маршрутизации транспортных потоков, управлении сервисами бронирования, образовательных платформах, управлении посетителями выставочных пространств. Несмотря на разнообразие приложений, существующие подходы не закрывают совокупности требований, характерной для рассматриваемой постановки: совместная оптимизация рекомендаций между пользователями, жёсткие ограничения вместимости в одной коллективной выдаче, ограниченность исторических данных, мягкий характер управляющего рычага.

Развитие рекомендательных систем, методов обучения с подкреплением и симуляционных подходов с использованием больших языковых моделей в 2023–2026 годах расширило инструментарий, применимый к задачам этого класса. Подходы на основе генеративных агентов и LLM-симуляторов допускают работу с малыми каталогами при отсутствии исторических данных взаимодействий; методы обучения политики с ограничениями дают формальный аппарат для учёта системных ограничений вместимости. Систематический синтез этих подходов в задаче формирования программы конференции составляет содержание настоящей работы.

\section*{Степень разработанности темы}
\addcontentsline{toc}{section}{Степень разработанности темы}

Задача в её обобщённой постановке исследовалась с разных сторон. В рамках теории игр аппарат игр с перегрузкой Розенталя [Rosenthal, 1973] и равновесия Уордропа [Wardrop, 1952] обеспечивают концептуальное описание ситуации децентрализованного выбора при взаимной зависимости пользователей. В рекомендательных системах подходы ранжирования с ограничениями Сингха и Йоахимса [Singh, Joachims, 2018] и калиброванных рекомендаций Стека [Steck, 2018] вводят аппарат учёта системных свойств выдачи. В обучении с подкреплением марковский процесс принятия решений с ограничениями [Altman, 1999] и его применения в производственных рекомендательных системах [Chen et al., 2019; Zheng et al., 2018] предоставляют операциональный механизм построения управляющей политики. Современные направления использования больших языковых моделей в роли симуляторов пользователей [Park et al., 2023; Zhang et al., 2024] и в роли ранжирующих компонентов [Hou et al., 2024] открывают возможность работы с малыми каталогами и качественной валидации моделей пользовательского поведения.

Вместе с тем, систематическое объединение указанных подходов применительно к задаче рассматриваемой структуры — формирование согласованных рекомендательных выдач для множества пользователей одновременно при жёстких ограничениях вместимости в одной коллективной выдаче — в открытой литературе отсутствует. Это определяет содержательное пространство для разработки методологического и алгоритмического подхода в работе.

\section*{Объект и предмет исследования}
\addcontentsline{toc}{section}{Объект и предмет исследования}

\textbf{Объект исследования} — процесс формирования персональных программ участников научно-практических IT-конференций при ограничениях вместимости параллельных сессий.

\textbf{Предмет исследования} — методы построения управляющей рекомендательной политики, согласующей индивидуальные предпочтения участников с системными ограничениями вместимости и обеспечивающей совместную оптимизацию рекомендаций между пользователями.

\section*{Цель работы}
\addcontentsline{toc}{section}{Цель работы}

Цель работы — разработать интеллектуальную систему поддержки формирования программы конференции, учитывающую ограничения вместимости параллельных сессий и согласованность рекомендаций между участниками.

\section*{Задачи исследования}
\addcontentsline{toc}{section}{Задачи исследования}

Для достижения поставленной цели в работе решаются следующие задачи:

\begin{enumerate}
  \item Провести анализ предметной области, систематизировать существующие концептуальные рамки и методы решения задач рассматриваемого класса, выделить пробелы в существующих подходах.
  \item Формализовать постановку задачи как марковский процесс принятия решений с ограничениями в структуре игры с перегрузкой; сформулировать функциональные и нефункциональные требования к разрабатываемому решению.
  \item Разработать гибридную симуляционную среду для обучения и оценки рекомендательных политик в условиях отсутствия исторических данных взаимодействий, объединяющую однократную генерацию синтетической пользовательской популяции на основе большой языковой модели и параметрический пошаговый симулятор.
  \item Реализовать набор сравниваемых рекомендательных политик различной природы — случайную, контентную на основе косинусной близости, эвристическую с учётом вместимости, основанную на алгоритме Maximal Marginal Relevance, обучаемую методом проксимальной оптимизации в постановке марковского процесса с ограничениями и опирающуюся на большую языковую модель — и провести их сравнение по совокупности метрик качества.
  \item Провести независимую валидацию политики, выбранной по результатам основного сравнения, в среде на основе автономных агентов с большими языковыми моделями; оценить разрыв между симуляцией и реальным поведением пользователей.
\end{enumerate}

\section*{Методы исследования}
\addcontentsline{toc}{section}{Методы исследования}

Работа опирается на анализ литературы по рекомендательным системам, теории игр с перегрузкой, обучению с подкреплением и стохастической оптимизации; математическое моделирование задачи как марковского процесса принятия решений с ограничениями; разработку симуляционной среды с использованием параметрических моделей пользовательского выбора и генеративных агентов на основе больших языковых моделей; обучение политики методами обучения с подкреплением; статистическое сравнение результатов нескольких рекомендательных политик по совокупности метрик качества.

\section*{Научная новизна}
\addcontentsline{toc}{section}{Научная новизна}

Научная новизна работы заключается в систематическом синтезе элементов нескольких теоретических рамок применительно к задаче рассматриваемого класса. Конкретно: сочетание формализации задачи как марковского процесса принятия решений с ограничениями со структурным обоснованием через теорию игр с перегрузкой, использование гибридной симуляционной среды с однократной генерацией пользовательской популяции на основе большой языковой модели для решения задач с малым каталогом ресурсов, двухслойная схема валидации с независимой проверкой в среде на основе автономных агентов. Указанное сочетание адресует совокупность требований, не закрываемую ни одним из существующих подходов в отдельности.

\section*{Практическая значимость}
\addcontentsline{toc}{section}{Практическая значимость}

Практическая значимость работы определяется применимостью разработанной системы к реальной задаче формирования программы конференций — на примере конференций Mobius и Heisenbug компании JUG. Эмпирическая часть выполняется на материале открытых программ конференций Mobius и Heisenbug 2023–2025 годов (около 30–40 докладов на конференцию с описаниями тематики, форматов и распределением по тайм-слотам и залам); параметры симуляционной среды калибруются по фактической вместимости залов и типичной численности участников. Разработанные симуляционная среда и набор сравниваемых политик могут быть переиспользованы для решения задач смежных доменов, входящих в рассматриваемый класс: маршрутизация транспортных потоков, бронирование сервисов с ограниченной вместимостью, управление учебной нагрузкой в образовательных платформах, управление потоками посетителей в зонах с пропускными ограничениями.

\section*{Положения, выносимые на защиту}
\addcontentsline{toc}{section}{Положения, выносимые на защиту}

\begin{enumerate}
  \item Постановка задачи формирования согласованных рекомендательных выдач при ограничениях вместимости как марковского процесса принятия решений с ограничениями в структуре игры с перегрузкой обеспечивает формальный аппарат для совместной оптимизации рекомендаций между пользователями при жёстких ограничениях вместимости в одной коллективной выдаче.
  \item Гибридная симуляционная среда, объединяющая однократную генерацию пользовательской популяции на основе большой языковой модели с параметрическим пошаговым симулятором, обеспечивает обучаемость рекомендательной политики при отсутствии реальных данных взаимодействий и контролируемой стоимости вычислений.
  \item Сравнительное исследование политик различной природы — случайной, контентной, эвристической с учётом вместимости, основанной на алгоритме MMR, обучаемой методом проксимальной оптимизации в постановке марковского процесса с ограничениями и опирающейся на большую языковую модель — позволяет систематически охарактеризовать компромиссы между соблюдением системных ограничений вместимости и пользовательской полезностью на материале рассматриваемой задачи.
  \item Двухслойная схема валидации с проверкой выбранной политики во второй симуляционной среде на основе автономных агентов с большими языковыми моделями позволяет оценить устойчивость результатов основного сравнения к смене модели пользовательского поведения и сформировать содержательную интерпретацию расхождений между двумя моделями симуляции.
\end{enumerate}

\section*{Апробация результатов}
\addcontentsline{toc}{section}{Апробация результатов}

Результаты работы получены и проверены на разработанной двухслойной симуляционной среде, объединяющей параметрический пошаговый симулятор и независимый слой валидации на основе автономных агентов с большими языковыми моделями. Подобная схема обеспечивает воспроизводимое сравнение методов и независимый контроль внешней валидности результатов. Постановка задачи и предлагаемый методологический подход обсуждались с научным руководителем работы и представителями команды разработки рекомендательного сервиса конференций Mobius и Heisenbug компании JUG, на материале которых выполняется эмпирическая часть работы.

\section*{Структура работы}
\addcontentsline{toc}{section}{Структура работы}

Работа состоит из введения, четырёх глав, заключения, списка использованных источников и приложений. В главе 1 проводится анализ предметной области, систематизируются концептуальные рамки и методы, идентифицируется пробел в существующих подходах. В главе 2 формализуется постановка задачи, формулируются требования к решению, описывается предлагаемый метод и определяется дизайн экспериментальной части работы. В главе 3 представляются технические решения по реализации симуляционной среды и сравниваемых политик, описывается ход проведённых экспериментов. В главе 4 приводятся результаты сравнения политик, ablation studies, проверка устойчивости и итоговая валидация в среде на основе автономных агентов; обсуждается разрыв между симуляцией и реальным поведением. В заключении формулируются основные результаты работы и направления дальнейших исследований.

\chapter{АНАЛИЗ ПРЕДМЕТНОЙ ОБЛАСТИ}


\section{Описание предметной области}

\textbf{Класс рассматриваемых задач.} В работе рассматривается класс задач распределения потока пользователей между ресурсами ограниченной вместимости при условии, что инструмент управления распределением — рекомендательная подсказка, не носящая принудительного характера. Цель управляющей стороны — максимизировать суммарную полезность распределения, не допуская при этом ни перегрузки отдельных ресурсов, ни их хронической недозагрузки.

Подобный класс задач возникает в широком спектре прикладных областей. К нему относятся задачи маршрутизации транспортных потоков, в которых навигационный сервис рекомендует водителям маршруты при ограниченной пропускной способности дорожной сети; задачи распределения мест в сервисах бронирования и заведениях обслуживания; задачи управления учебной нагрузкой при ограничениях на численность учебных групп; задачи управления потоками посетителей в зонах с ограниченной пропускной способностью (выставочные пространства, парки, конференции); задачи управления нагрузкой в распределённых вычислительных системах. Несмотря на различия в предметной специфике, перечисленные задачи объединены общей структурой: наличие неоднородного спроса со стороны автономных пользователей, ограниченность ресурсов, отсутствие возможности принудительного назначения и наличие управляющего рычага в виде рекомендательной выдачи.

\textbf{Действующие стороны.} В постановке задачи выделяются три стороны.

Пользователи — субъекты, самостоятельно принимающие решение о выборе ресурса. Множество пользователей характеризуется неоднородностью предпочтений и автономностью поведения.

Ресурсы — объекты выбора. Каждый ресурс характеризуется ограничением по числу пользователей, превышение которого приводит к деградации качества обслуживания: отказам в доступе, ухудшению опыта, снижению эффективности.

Организатор — управляющая сторона, заинтересованная в качестве функционирования системы как целого. Не обладает полномочиями принудительного назначения; влияет на распределение исключительно через формирование рекомендательных подсказок — ранжированных списков, предъявляемых пользователю в качестве предпочтительных вариантов.

\textbf{Особенности задачи.} Управляющий механизм носит вероятностный характер: пользователь не обязан следовать рекомендации. Это отличает задачу от классических задач оптимального назначения, в которых решение управляющей стороны эквивалентно фактическому распределению.

Существенной особенностью является взаимосвязанность рекомендаций. Поскольку пользователи получают подсказки одновременно, рекомендация одному пользователю не может рассматриваться изолированно от рекомендаций, выдаваемых другим. При независимом формировании рекомендаций по индивидуальной релевантности множество пользователей получит в верхних позициях одни и те же объективно наиболее популярные ресурсы, что приведёт к ожидаемому совокупному потоку, превышающему вместимость. Следовательно, оптимальная индивидуальная подсказка для пользователя есть функция от распределения подсказок, выдаваемых остальным пользователям. Задача переходит в плоскость совместной оптимизации.

\textbf{Полезность и противоречие интересов.} Каждое сопоставление пользователя и ресурса характеризуется числовой оценкой полезности — степени соответствия ресурса предпочтениям пользователя с учётом контекстных факторов (времени, заполненности, удобства). Целевая функция организатора — суммарная полезность по всем парам «пользователь–ресурс», подлежащая максимизации при соблюдении ограничений вместимости.

Особенностью класса задач является системное противоречие между индивидуальной и совокупной оптимальностью. Индивидуальный выбор пользователей в пользу наиболее предпочтительного ресурса при отсутствии координации приводит к коллективному результату, заведомо уступающему оптимальному. Управляющая сторона решает задачу нивелирования данного противоречия: смещения индивидуального поведения в направлении системного оптимума при сохранении автономности пользователей.

\textbf{Конкретная инстанция предметной области.} В данной работе общий класс задач рассматривается на материале формирования программы научно-практических IT-конференций. Выбор инстанции обусловлен явной выраженностью характерных черт класса: значительной неоднородностью предпочтений участников, существенным разбросом вместимости параллельных сессий, отсутствием полной информации о пользователях, а также наличием доступа к экспериментальной площадке. Эмпирическая часть работы выполняется на данных конференций Mobius и Heisenbug компании JUG.

\section{Сущность решаемой проблемы}

\textbf{Деградация системы при отсутствии управления.} В отсутствие координирующего вмешательства распределение пользователей по ресурсам формируется как агрегат независимых индивидуальных выборов. Поскольку индивидуальные предпочтения неоднородны, но имеют общие пики (наиболее популярные позиции), результирующее распределение характеризуется одновременной перегрузкой части ресурсов и хронической недозагрузкой другой части. В терминах теории игр такое состояние — равновесие эгоистичных агентов, отстоящее от системного оптимума на величину, измеряемую показателем «цена анархии» (Price of Anarchy).

\textbf{Структура потерь.} Потери при подобном состоянии затрагивают всех участников взаимодействия. Конкретный их характер варьируется в зависимости от природы домена; ниже описаны категории, имеющие универсальный характер для всего рассматриваемого класса задач.

Пользователи, обращающиеся к перегруженным ресурсам, сталкиваются с отказами в обслуживании, увеличением времени ожидания, снижением качества обслуживания, отсутствием равноценных альтернатив в момент возникновения конфликта. Конкретное стоимостное выражение этих издержек определяется доменом: прямые денежные затраты, не компенсируемые получаемой полезностью; потеря времени; упущенные возможности; ухудшение качества предоставляемой услуги.

Управляющая сторона несёт потери, связанные с неэффективным использованием инфраструктурного капитала: затраты на обеспечение недозагруженных ресурсов не компенсируются объёмом фактического использования. Снижение интегрального качества обслуживания одновременно отражается на показателях лояльности пользователей и продолжительности их участия в системе. При наличии у управляющей стороны коммерческих интересов, опосредованных распределением пользователей — например, интересов внешних партнёров, обусловленных охватом отдельных сегментов системы, — неравномерность распределения дополнительно снижает соответствующую коммерческую ценность.

В доменах, где роль поставщика ресурсов выделена как самостоятельная, поставщики недозагруженных ресурсов сталкиваются с обесцениванием своих ресурсов как канала распространения предоставляемого через них содержания. Это снижает мотивацию поставщиков к продолжению участия в системе и в долгосрочной перспективе ведёт к снижению разнообразия и качества предлагаемых ресурсов.

\textbf{Несостоятельность прямого применения стандартных подходов.} Прямое применение классической рекомендательной системы не решает поставленной задачи. Классическая рекомендательная система оптимизирует индивидуальную полезность пользователя на основании его предпочтений, не учитывая ограничений вместимости и взаимосвязанности рекомендаций. При одновременной выдаче рекомендаций множеству пользователей подобная система воспроизводит одинаково популярные позиции в верхних строках выдачи, тем самым усиливая, а не ослабляя нагрузку на наиболее востребованные ресурсы.

Альтернативная стратегия централизованного принудительного назначения также неприменима. С одной стороны, она входит в противоречие с автономностью пользователей и обесценивает само участие, основанное на свободе выбора. С другой стороны, формальное назначение не гарантирует фактического следования предписанию: реальное поведение пользователей сохраняет вероятностный характер вне зависимости от формы предъявления рекомендации.

\textbf{Природа сложности задачи.} Сложность задачи определяется одновременным выполнением трёх требований, каждое из которых разрешимо в отдельности, а в совокупности образующих противоречивую систему:

\begin{enumerate}
  \item Соблюдение индивидуальной релевантности: каждая выдаваемая рекомендация должна оставаться приемлемой для пользователя с точки зрения его предпочтений; в противном случае пользователь игнорирует подсказку.
  \item Соблюдение ограничений вместимости: суммарный ожидаемый поток пользователей на каждый ресурс не должен превышать его вместимости.
  \item Согласованность совокупной выдачи: множество персональных рекомендаций образует когерентное распределение; рекомендация одного пользователя учитывает рекомендации, выдаваемые другим.
\end{enumerate}

Третье требование принципиально отличает задачу от классической постановки рекомендательных систем, в которой пользователи рассматриваются как независимые единицы анализа. В рассматриваемом классе задач пользовательские рекомендации становятся взаимозависимыми, что переводит задачу из множества независимых ранжирований в задачу совместной оптимизации с ограничениями. Объём пространства решений и вычислительная сложность возрастают при этом нелинейно с ростом числа пользователей и ресурсов.

Указанная взаимосвязанность представляет основной источник сложности задачи и обусловливает необходимость отдельного методологического подхода, не сводимого к прямому применению существующих методов классических рекомендательных систем.

\section{Обзор существующих подходов и решений}

Описанный в подразделах 1.1–1.2 класс задач допускает постановку в нескольких самостоятельных формальных рамках, каждая из которых предлагает собственный понятийный аппарат, классы алгоритмов и методические инструменты. Ниже последовательно рассматриваются пять рамок: исследование операций, теория игр, рекомендательные системы с ограничениями, обучение с подкреплением и стохастическая оптимизация. Рамки не являются взаимоисключающими: в практической работе они пересекаются и комбинируются. Конкретный методологический синтез, выбранный в работе, представлен в главе 2.

\subsection{Постановка в рамках исследования операций}

Исследование операций (Operations Research) рассматривает задачу как комбинаторную задачу распределения дискретных пользователей по ограниченным ресурсам с целью максимизации суммарной полезности. Каноническая формулировка — обобщённая задача о назначениях (Generalized Assignment Problem, GAP), являющаяся классическим объектом изучения комбинаторной оптимизации.

Формальная постановка имеет следующий вид. Для множества пользователей $I = \{1, \ldots, n\}$, множества ресурсов $J = \{1, \ldots, m\}$, функции полезности $u_{ij}$ и вместимостей $c_j$ требуется найти бинарные индикаторы назначения $x_{ij} \in \{0, 1\}$, максимизирующие суммарную полезность $\sum_{ij} u_{ij} x_{ij}$ при условии однократного назначения каждого пользователя ($\sum_j x_{ij} = 1$ для всех $i$) и соблюдения вместимостей ($\sum_i x_{ij} \leq c_j$ для всех $j$). В случае, когда пользователю необходимо последовательно использовать несколько ресурсов в различные временные интервалы, постановка расширяется: ограничение однократного назначения заменяется условием, согласно которому в каждый момент времени за пользователем закреплено не более одного ресурса. Подобная расширенная формулировка относится к классу многопериодных задач о назначениях и пересекается с задачами составления расписаний (timetabling), для которых разработан самостоятельный методический аппарат.

С алгоритмической точки зрения GAP относится к классу NP-полных задач, что в общем случае исключает существование точного решения в полиномиальное время. Практически применяемые методы включают методы ветвей и границ для точного решения экземпляров умеренной размерности, лагранжеву релаксацию для построения оценок и приближённых решений [Fisher, Jaikumar, Van Wassenhove, 1986], а также эвристики на основе жадных алгоритмов и локального поиска. Систематический обзор алгоритмов представлен в работе Cattrysse и Van Wassenhove [1992]. Каноническим источником, охватывающим как задачи о рюкзаке, так и задачи о назначениях, является монография Martello и Toth [1990].

В прикладном плане аппарат исследования операций применяется в задачах распределения трудовых ресурсов (назначения работников на задачи с учётом нагрузки), в задачах распределения учащихся по образовательным группам, в задачах планирования производственных мощностей. В каждом из этих приложений решение задачи трактуется как итоговое распределение, подлежащее исполнению.

Рамка исследования операций даёт строгий формализм и развитый алгоритмический арсенал. Применительно к рассматриваемой задаче ограничивается допущением о детерминированном исполнении решения: в постановке GAP $x_{ij} = 1$ означает фактическое назначение, тогда как в работе решение управляющей стороны имеет статус рекомендации с вероятностным откликом. Прямое применение GAP даёт решения, оптимальные в детерминированной модели, но не учитывающие отклика пользователей.

\subsection{Постановка в рамках теории игр}

Теория игр рассматривает задачу как ситуацию совместного принятия решений автономными агентами, каждый из которых преследует собственную выгоду, а коллективное распределение возникает как результат отдельных индивидуальных выборов. Естественной формальной рамкой для рассматриваемого класса задач является игра с перегрузкой (congestion game) — класс игр, в которых полезность игрока на ресурсе зависит от числа других игроков, выбравших тот же ресурс.

Каноническая формулировка дана Розенталем [Rosenthal, 1973]. Задаётся множество игроков $N$, множество ресурсов $E$ и функция издержек $c_e: \mathbb{N} \to \mathbb{R}$ для каждого ресурса $e$, зависящая от числа игроков, выбравших этот ресурс. Стратегия игрока $i$ — подмножество ресурсов $s_i \subseteq E$. Издержки игрока определяются как $\sum_{e \in s_i} c_e(n_e(s))$, где $n_e(s)$ — число игроков, использующих ресурс $e$ в совместной стратегии $s$. Розенталем установлено, что любая такая игра обладает равновесием Нэша в чистых стратегиях и относится к классу точных потенциальных игр: существует единая функция потенциала $\Phi(s) = \sum_e \sum_{k=1}^{n_e(s)} c_e(k)$, монотонно изменяющаяся при любом одностороннем улучшении стратегии игрока. Концепция потенциальных игр впоследствии обобщена в работе Мондерера и Шэпли [Monderer, Shapley, 1996].

В задачах с большим числом однотипных пользователей, в которых индивидуальный вклад одного пользователя пренебрежимо мал по отношению к общему потоку, применяется непрерывная постановка, восходящая к работе Уордропа [Wardrop, 1952] о моделировании транспортных потоков. Согласно первому принципу Уордропа, в состоянии равновесия все используемые маршруты между парой источник–назначение имеют одинаковую стоимость, не превышающую стоимости любого неиспользуемого маршрута. Второй принцип Уордропа описывает системно-оптимальное распределение, минимизирующее суммарные издержки. Разрыв между равновесием эгоистичных пользователей и системным оптимумом является центральной характеристикой задачи.

Количественная мера данного разрыва — показатель «цена анархии» (Price of Anarchy), введённый в работе Кутсупиаса и Пападимитриу [Koutsoupias, Papadimitriou, 1999] и определяемый как отношение стоимости наихудшего равновесия к стоимости системного оптимума. Для задач эгоистичной маршрутизации с линейными функциями задержек верхняя граница цены анархии составляет 4/3 [Roughgarden, Tardos, 2002]. Иллюстрацией непредсказуемости системного поведения при индивидуальном выборе служит парадокс Брэсса [Braess, 1968], показывающий, что добавление дополнительного маршрута в сеть может ухудшить равновесное распределение и увеличить издержки всех пользователей. Систематическое изложение результатов теории эгоистичной маршрутизации представлено в монографии Раугардена [Roughgarden, 2005].

В прикладном плане теоретико-игровая постановка наиболее активно используется в задачах маршрутизации транспортных потоков, в которых навигационные сервисы реализуют рассматриваемый сценарий непосредственно: каждый пользователь стремится минимизировать собственное время в пути, а заторы возникают как естественное следствие отсутствия координации между пользователями.

Теоретико-игровая рамка естественно соответствует природе задачи (автономные пользователи, децентрализованные решения, агрегатное распределение); аппарат равновесий и цены анархии даёт теоретическую базу для анализа устойчивых состояний системы. Применительно к управляющей задаче эта рамка остаётся описательной — она отвечает на вопрос, какое распределение установится, но не даёт алгоритма построения смещающего воздействия. Переход к проектированию управляющих воздействий относится к смежным направлениям (mechanism design, information design), не предоставляющим прямых вычислительных алгоритмов для рекомендательных задач.

\subsection{Постановка в рамках рекомендательных систем с ограничениями}

Подход рекомендательных систем рассматривает задачу как формирование индивидуального ранжированного списка для каждого пользователя на основании предсказанной полезности, при этом классическая постановка дополняется системными ограничениями на агрегированное распределение выдач. В отличие от классических задач ранжирования, оптимизирующих исключительно индивидуальную релевантность, постановки с ограничениями вводят дополнительные критерии, относящиеся к свойствам совокупной выдачи: соблюдение распределительных пропорций, сохранение справедливости по отношению к поставщикам ресурсов, соблюдение системных ограничений вместимости.

Ранжирование с ограничениями (constrained ranking) — формальная постановка, систематически разработанная в работах Сингха и Йоахимса [Singh, Joachims, 2018; 2019]. Политика ранжирования $\pi$, отображающая пользователя в распределение по позициям выдачи, оптимизируется с целью максимизации ожидаемой пользовательской полезности при дополнительных линейных ограничениях на агрегированную экспозицию ресурсов:

\[
\max_\pi \mathbb{E}_{u} \left[ \sum_i \pi_i(u) r_{ui} \right] \quad \text{при} \quad \sum_{u} \pi_i(u) \in [c_i^{\min}, c_i^{\max}] \ \forall i,
\]

где $r_{ui}$ — предсказанная релевантность ресурса $i$ для пользователя $u$, а $c_i^{\min}, c_i^{\max}$ — нижние и верхние границы совокупной экспозиции ресурса $i$. Авторами показано, что задача допускает эффективное решение методами линейного программирования и применима как для оптимизации справедливости экспозиции, так и для соблюдения системных ограничений на распределение.

Калиброванные рекомендации (calibrated recommendations), предложенные Стеком [Steck, 2018], представляют альтернативную перспективу: для каждого пользователя распределение рекомендованных ресурсов по категориям должно соответствовать его историческому профилю. Постановка опирается на минимизацию расхождения Кульбака–Лейблера между распределением рекомендаций и пользовательским профилем как дополнительного слагаемого в целевой функции. Хотя задача калибровки сформулирована на уровне индивидуальной выдачи, аппарат расхождения распределений применяется для измерения системных свойств совокупной выдачи, в том числе равномерности заполнения ресурсов.

Концепция справедливости со стороны поставщиков (provider-side fairness) разработана в направлении исследований многосторонних рекомендательных систем [Abdollahpouri, Burke, Mobasher, 2020]. В её рамках в качестве заинтересованных сторон признаются не только пользователи, но и поставщики ресурсов, ожидающие равномерного или пропорционального распределения внимания. Классические работы данного направления включают Биегу, Гуммади и Вайкума [Biega, Gummadi, Weikum, 2018] о равенстве внимания (equity of attention), Целике и соавторов [Zehlike et al., 2017] о справедливом ранжировании FA*IR, а также позднейшие обзоры [Naghiaei, Rahmani, Deldjoo, 2022], систематизирующие подходы к учёту провайдерских интересов в системах ранжирования.

В прикладном плане рекомендательные системы с ограничениями применяются в стриминговых сервисах для управления экспозицией содержания (Netflix, Spotify), в торговых площадках для распределения внимания между продавцами, в задачах формирования ленты социальных сетей для управления разнообразием выдачи.

Рекомендательная рамка даёт развитый аппарат операциональных метрик (NDCG, MRR, ILD, метрики калибровки и справедливости экспозиции) и алгоритмы переранжирования, применимые поверх базовой модели релевантности. Главное ограничение для рассматриваемой задачи — формулировка ограничений как математического ожидания по распределению пользователей (выравнивание «в среднем» по большому потоку выдач), тогда как в работе требуется соблюдение вместимости в каждой конкретной коллективной выдаче. Адаптация классических методов под одноразовый совместный режим требует существенных модификаций.

\subsection{Постановка в рамках обучения с подкреплением}

Подход обучения с подкреплением (reinforcement learning, RL) рассматривает задачу как последовательный процесс взаимодействия между управляющей политикой и пользователями, в котором политика формирует рекомендательные подсказки, наблюдает отклик и адаптирует своё поведение на основании накопленного опыта взаимодействия. В отличие от подходов, оптимизирующих заранее зафиксированную целевую функцию на статической модели предпочтений, рамка обучения с подкреплением позволяет обучать политику непосредственно из взаимодействий с пользователями или их моделью, учитывая отложенные эффекты решений и динамику пользовательского поведения.

Каноническая формализация — марковский процесс принятия решений (Markov Decision Process, MDP). Управляющий агент в каждом шаге наблюдает состояние $s_t$ (контекст пользователя, история взаимодействия, состояние системы), выбирает действие $a_t$ (рекомендацию или ранжированный список), получает мгновенную награду $r_t$ (показатель полезности — клик, удовлетворение, продолжительность взаимодействия) и переходит в новое состояние $s_{t+1}$ согласно вероятности перехода $P(s_{t+1} | s_t, a_t)$. Целью является нахождение политики $\pi$, максимизирующей дисконтированную суммарную награду $\mathbb{E}_\pi \left[ \sum_{t=0}^\infty \gamma^t r_t \right]$.

Постановка задачи с ограничениями — марковский процесс с ограничениями (Constrained MDP), систематически разработанная в монографии Альтмана [Altman, 1999]. К стандартной MDP добавляется вектор стоимостных функций $c_1, \ldots, c_k$; ограничения принимают вид $\mathbb{E}_\pi \left[ \sum_t \gamma^t c_{i,t} \right] \leq d_i$. Альтманом доказано существование оптимальной стационарной стохастической политики и сводимость задачи к линейному программированию на пространстве распределений по парам «состояние–действие». В терминах рассматриваемого класса задач функции $c_i$ соответствуют ограничениям вместимости ресурсов, а пороговые значения $d_i$ — допустимой загрузке.

В качестве промежуточной рамки между классическим обучением с подкреплением и комбинаторной оптимизацией используются многорукие бандиты с ограничениями ресурсов (bandits with knapsacks), формализованные в работах Баданидиюру и соавторов [Badanidiyuru et al., 2013]. В этой постановке каждый раунд рассматривается как одно взаимодействие с пользователем без учёта долгосрочных переходов состояний, а глобальные ограничения формулируются как ограничения на суммарное потребление ресурсов. Линейная контекстная версия задачи [Agrawal, Devanur, 2016] обеспечивает алгоритмы с гарантированной верхней оценкой сожаления порядка $\tilde{O}(\sqrt{T})$.

Применение обучения с подкреплением в рекомендательных системах развивается с середины 2000-х; современный этап представлен работами о глубоком обучении с подкреплением для рекомендаций [Zheng et al., 2018], производственным применением метода REINFORCE с коррекцией смещения вне политики на крупной видеоплатформе [Chen et al., 2019] и систематическими обзорами направления [Afsar et al., 2022]. Постановки с ограничениями получили развитие в работе [Zheng et al., 2022] и в технических отчётах индустриальных лабораторий [Netflix Tech Blog, 2022], демонстрирующих жизнеспособность подхода в промышленных условиях. Большинство открытой литературы опирается на массивные объёмы накопленных данных взаимодействий, обеспечивающие надёжное оффлайн-обучение и валидацию методами оценки вне политики.

Рамка обучения с подкреплением учитывает динамику взаимодействия и отложенные эффекты решений; марковский процесс с ограничениями обеспечивает формальную базу учёта системных ограничений. Для рассматриваемой задачи характерны два ограничения. Первое — высокие требования к объёму данных: производственные применения опираются на миллионы взаимодействий, что недоступно при малом каталоге; следствие — необходимость симулятора пользовательского поведения и перенос части валидационной нагрузки в плоскость качества симулятора. Второе — формулировка ограничений как математического ожидания по политике, означающая, как и в рекомендательной рамке, выравнивание «в среднем», а не соблюдение требований в каждой коллективной выдаче.

\subsection{Постановка в рамках стохастической оптимизации}

Подход стохастической оптимизации рассматривает задачу с явным учётом случайности пользовательского поведения. Каноническими формулировками являются двухэтапное стохастическое программирование [Birge, Louveaux, 2011], в котором решение принимается в условиях неопределённости и корректируется после её раскрытия; вероятностные ограничения [Charnes, Cooper, 1959], формулируемые как требования на вероятность нарушения; робастная оптимизация [Ben-Tal, El Ghaoui, Nemirovski, 2009], дающая гарантии работоспособности по наихудшему сценарию. Прикладным образцом служит область управления доходностью (revenue management) [Talluri, van Ryzin, 2004], структурно близкая к рассматриваемой задаче.

В настоящей работе стохастическая рамка не выступает основным методологическим аппаратом; элементы вероятностного описания используются в формализации задачи (раздел 2.1), но операциональным механизмом построения политики выбран марковский процесс принятия решений с ограничениями. Высокая вычислительная сложность точных стохастических постановок и необходимость в данных для оценивания распределения спроса ограничивают их прямое применение в условиях малого каталога ресурсов и отсутствия исторических данных взаимодействий.

\section{Обзор технологий и методов}

Если в подразделе 1.3 систематизированы концептуальные рамки постановки задачи, то здесь приводится обзор конкретных алгоритмов, моделей и подходов, разработанных в литературе и индустрии. Цель — определить инструментарий, доступный для непосредственного использования или адаптации, и оценить его применимость с учётом специфики: малый каталог ресурсов, отсутствие исторических данных взаимодействий, требование соблюдения ограничений в одноразовой коллективной выдаче. Обзор проводится по четырём направлениям: алгоритмы переранжирования с системными ограничениями (1.4.1), методы обучения с подкреплением для рекомендательных задач (1.4.2), параметрические симуляционные среды (1.4.3), подходы на основе больших языковых моделей и автономных агентов (1.4.4).

\subsection{Алгоритмы переранжирования с системными ограничениями}

Переранжирование — наиболее доступный класс инструментов учёта системных ограничений в рекомендательной выдаче: к базовому ранжированию по индивидуальной релевантности применяется преобразование, обеспечивающее заданные свойства совокупного результата (разнообразие, справедливость, баланс категорий, ограничения вместимости). Базовая модель не переобучается, что упрощает внедрение в существующие системы.

Канонические алгоритмы класса включают Maximal Marginal Relevance [Carbonell, Goldstein, 1998] — жадный метод балансирования релевантности и разнообразия через параметр $\lambda$; детерминантные точечные процессы [Kulesza, Taskar, 2012], выбирающие подмножество с вероятностью, пропорциональной определителю подматрицы ядра, что штрафует похожие элементы; алгоритм FA*IR [Zehlike et al., 2017] для справедливого формирования top-k выдачи с гарантиями представления защищённой группы; алгоритм калиброванных рекомендаций [Steck, 2018], минимизирующий расхождение Кульбака–Лейблера между распределением выдачи по категориям и заданным эталонным распределением. Последний наиболее близок к задачам соблюдения распределительных ограничений на уровне коллективной выдачи: при надлежащем определении эталона аппарат расхождения распределений непосредственно применим к контролю заполненности ресурсов.

Общее ограничение алгоритмов переранжирования применительно к рассматриваемой задаче — их статический характер. Алгоритм оптимизирует выдачу при заданных предположениях о пользовательском поведении, но не учитывает вероятностного отклика и взаимной зависимости рекомендаций между пользователями; распространение на коллективную оптимизацию требует существенных обобщений, выходящих за рамки оригинальных постановок.

\subsection{Методы обучения с подкреплением для рекомендательных задач}

Методы обучения с подкреплением, применяемые в рекомендательных системах, делятся на три семейства: основанные на функции ценности (Q-обучение [Watkins, Dayan, 1992] и его глубокая версия DQN [Mnih et al., 2015], применённая в DRN [Zheng et al., 2018] для рекомендации новостей); основанные на политике (REINFORCE [Williams, 1992]); гибридные методы «актор-критик» (A2C/A3C [Mnih et al., 2016] и PPO [Schulman et al., 2017]). Семейство «актор-критик» естественно интегрируется с задачами с ограничениями через двойственную оптимизацию множителей Лагранжа [Zheng et al., 2022].

Производственные применения опираются на накопленные журналы взаимодействий и требуют методов компенсации смещения между обучающей и целевой политиками — выборки по значимости (importance sampling) и её устойчивых вариантов, а также методов дважды робастной оценки. Работа Чена и соавторов [Chen et al., 2019] предложила Top-K коррекцию вне политики для метода REINFORCE, обеспечив развёртывание обучаемой политики в крупной производственной системе.

В рассматриваемой задаче прямое применение обучения с подкреплением осложнено отсутствием накопленных журналов взаимодействий: малый объём исторических данных делает методы оценки вне политики ненадёжными. Альтернативой служит обучение и оценка политик в симуляционной среде (подраздел 1.4.3).

\subsection{Параметрические симуляционные среды для рекомендательных систем}

При недоступности достаточных объёмов накопленных данных взаимодействий и невозможности масштабных онлайн-экспериментов основным альтернативным источником данных служит симуляционная среда, моделирующая поведение пользователей в ответ на рекомендательные политики. Параметрический класс симуляторов задаёт поведение через явно специфицированные вероятностные функции выбора; альтернативный класс на основе больших языковых моделей рассматривается в подразделе 1.4.4.

Каноническими представителями класса являются RecSim [Ie et al., 2019] (Google; модель пользователя с латентным состоянием, моделями выбора и интеграцией с библиотеками обучения с подкреплением); Sim4Rec [Sber AI Lab, 2022] (открытая разработка ПАО «Сбер», ориентированная на автономные эксперименты с поддержкой синтетической генерации пользователей и моделей выбора, имеющая документацию на русском языке); KuaiSim [Zhao et al., 2023] (среда исследовательской группы Kuaishou с поведенческой моделью, обученной на журналах производственной системы) с сопровождающим набором полностью наблюдаемых данных KuaiRec [Gao et al., 2022].

Параметрические симуляционные среды применимы к рассматриваемой задаче, однако сопровождаются известными методологическими оговорками: разрыв между симулируемым и реальным поведением пользователей (sim-to-real gap) ставит вопрос внешней валидности; сам симулятор содержит содержательные предположения о поведении, требующие явного обоснования; при использовании среды как основной площадки сравнения методов необходима независимая проверка устойчивости результатов к выбору параметров.

\subsection{Подходы на основе больших языковых моделей и автономных агентов}

В 2023–2026 годах сформировался самостоятельный класс подходов, основанных на использовании больших языковых моделей (large language models, LLM) и автономных агентов на их базе. В отличие от параметрических методов, агентный подход возлагает функцию принятия решений на экземпляр LLM, получающей описание персоны пользователя в естественном языке. Большая языковая модель может выступать также в роли ранжирующего компонента, генератора синтетических данных, инструмента оценки качества или элемента контура обучения.

В роли симуляторов пользователей идея восходит к работе о генеративных агентах [Park et al., 2023] и развивается в серии работ для рекомендательного домена — Agent4Rec [Zhang et al., 2024], AgentCF, RecAgent — а также в крупномасштабных средах OASIS [Yang et al., 2024] и AgentSociety [Piao et al., 2025], моделирующих миллионы агентов в социальных и городских контекстах.

В роли ранжирующего компонента современные LLM способны выполнять ранжирование без дополнительного обучения (zero-shot) [Hou et al., 2024], интегрируясь поверх существующих систем поиска по эмбеддингам. Семейство методов LLMRec, P5, GPT4Rec и ChatRec, а также свежие подходы BIGRec [Bao et al., 2024] и LLaRA [Liao et al., 2024] исследуют разные способы использования LLM в задачах рекомендации — от прямой генерации списков до гибридных архитектур, объединяющих эмбеддинги традиционных рекомендательных моделей с языковыми моделями. Агентные рекомендательные системы развивают подход через паттерны ReAct [Yao et al., 2022] и Reflexion.

В задачах генерации синтетических данных и оценки качества LLM применяются для создания пользовательских профилей и сценариев поведения, что особенно ценно для задач с малым каталогом ресурсов. Парадигма LLM-as-judge [Zheng et al., 2023] обеспечивает ускоренную автономную оценку рекомендательных политик без дорогостоящих пользовательских исследований.

Гибридные схемы комбинируют преимущества обоих классов: однократная LLM-генерация персон с последующим использованием в параметрическом симуляторе существенно снижает совокупную вычислительную стоимость при сохранении содержательного разнообразия популяции. Двухслойная валидация — обучение в параметрическом симуляторе и независимая проверка в LLM-агентной среде — согласует требования эффективности и достоверности результатов. Главные ограничения подхода — вычислительная стоимость прямого использования LLM на каждом шаге и вопрос соответствия поведения LLM-агентов поведению реальных пользователей, активно обсуждаемый в свежей литературе.

\section{Анализ ограничений существующих решений}

Систематизация рассмотренных в подразделах 1.3–1.4 концептуальных рамок, методов и инструментов позволяет выделить ключевые ограничения существующих подходов. Анализ проводится по нескольким независимым осям, каждая из которых отражает важное свойство задачи, недостаточно покрываемое существующей литературой.

\textbf{Ось 1. Совместная оптимизация рекомендаций между пользователями.} Большинство рамок рассматривают задачу на уровне индивидуального пользователя; постановка, в которой множество пользователей одновременно требуют согласованной выдачи, представлена лишь в исследовании операций и в специализированных постановках управления доходностью, причём первое не учитывает вероятностного отклика.

\textbf{Ось 2. Жёсткие ограничения вместимости в одноразовой выдаче.} Подходы рекомендательных систем с ограничениями и постановки Constrained MDP формулируют ограничения как математическое ожидание по распределению пользователей или по политике, что обеспечивает выравнивание «в среднем» по большому потоку выдач, но не соблюдение в каждой конкретной коллективной выдаче, требуемое в рассматриваемой постановке.

\textbf{Ось 3. Ограниченность исторических данных.} Производственные применения опираются на миллионы накопленных взаимодействий; в задачах с малым каталогом и ограниченным числом пользователей подобный объём данных принципиально недоступен, что исключает прямое применение методов обучения с подкреплением и методов оценки вне политики.

\textbf{Ось 4. Природа управляющего рычага.} Классические постановки покрывают полярные случаи — либо детерминированное исполнение, либо отсутствие управляющего воздействия. Промежуточный случай мягкой рекомендации, которой пользователь может не следовать, формализован слабее; смежные направления (mechanism design, information design) развиты для теоретического анализа и не предоставляют прямых вычислительных алгоритмов.

\textbf{Ось 5. Симуляционная валидация при ограниченных данных.} Параметрические симуляторы (RecSim, Sim4Rec, KuaiSim) ориентированы на сценарии с большими каталогами и параметрическими моделями предпочтений, что несёт риск упрощения при малом каталоге. LLM-агентные симуляторы (подраздел 1.4.4) дают содержательную глубину, но сопровождаются высокой вычислительной стоимостью и открытым вопросом соответствия реальному поведению пользователей.

\textbf{Идентифицированный пробел.} Объединение перечисленных осей формирует область, систематически не покрытую существующими решениями: задача формирования согласованных рекомендательных выдач для множества пользователей одновременно, с жёсткими ограничениями вместимости в каждой одноразовой выдаче, при ограниченных исторических данных, в условиях мягкого характера управляющего рычага и при необходимости опоры на симуляционную валидацию. Привлечение средств LLM-агентного аппарата (подраздел 1.4.4) существенно расширяет возможности работы с осями 3 и 5 (ограниченность данных и симуляционная валидация), однако само по себе не закрывает оси 1 и 2 (совместная оптимизация и жёсткие ограничения вместимости в одноразовой выдаче). Ни одна из рассмотренных концептуальных рамок и ни один из существующих программных инструментов — включая инструменты, основанные на больших языковых моделях, — не предоставляет готового решения для указанного пересечения требований; работы, систематически адресующие все пять осей одновременно, в открытой литературе на момент проведения исследования не обнаружены. Этот пробел определяет содержательное пространство для разработки методологического и алгоритмического подхода в работе; конкретное наполнение пространства представлено в главе 2.

\section{Выводы по главе}

Рассматриваемый класс задач — распределение автономных пользователей по ресурсам ограниченной вместимости через рекомендательную подсказку, не носящую принудительного характера, — возникает в широком спектре прикладных областей и образует устойчивую предметную постановку. Конкретной инстанцией работы выбрано формирование программы научно-практических IT-конференций.

Сложность задачи определяется одновременным выполнением трёх требований: соблюдения индивидуальной релевантности рекомендаций, соблюдения ограничений вместимости ресурсов и согласованности совокупной выдачи между пользователями. Прямое применение классических рекомендательных систем эту задачу не закрывает, поскольку независимая оптимизация индивидуальной полезности усиливает, а не ослабляет нагрузку на наиболее востребованные ресурсы.

Среди существующих концептуальных рамок ни одна не покрывает совокупности характерных для задачи требований. Аппарат исследования операций ограничен предположением о детерминированном исполнении решения; классическая теория игр носит описательный, а не предписывающий характер; постановки рекомендательных систем с ограничениями формулируют ограничения как амортизированные ожидания, а не как требования к одной выдаче; методы обучения с подкреплением требуют объёмов данных, недоступных в условиях малого каталога; стохастическая оптимизация не предоставляет операциональной схемы построения управляющей политики. Аналогичный вывод получен для рассмотренных алгоритмов и симуляционных сред — ни один из готовых инструментов не закрывает специфики задачи без существенной адаптации.

Идентифицированный пробел, описанный в подразделе 1.5, определяет содержательное пространство работы. В главе 2 формализуется постановка задачи и описывается предлагаемый методологический подход, опирающийся на синтез аппарата марковского процесса принятия решений с ограничениями и теории игр с перегрузкой.

\chapter{МЕТОДЫ И ПОСТАНОВКА ЗАДАЧИ}

\section{Формализация задачи}

В подразделе 1.5 идентифицирован пробел в существующих решениях, относящийся к задаче формирования согласованных рекомендательных выдач для множества пользователей одновременно при жёстких ограничениях вместимости. Настоящий подраздел представляет формальную постановку задачи, опирающуюся на синтез двух концептуальных рамок, рассмотренных в подразделах 1.3.2 и 1.3.4: теории игр с перегрузкой (congestion games) и марковского процесса принятия решений с ограничениями (Constrained MDP). Первая рамка обеспечивает структурное описание ситуации, в которой коллективное распределение пользователей формируется как результат децентрализованных индивидуальных выборов под влиянием информационного сигнала; вторая — операциональный аппарат для построения управляющей политики, обучаемой методами обучения с подкреплением.

\subsection{Объекты и переменные задачи}

Рассматривается дискретное множество ресурсов, обозначаемых индексом из множества $J = \{1, \ldots, m\}$, и дискретное множество пользователей с индексами из множества $I = \{1, \ldots, n\}$. Каждый ресурс $j$ характеризуется вместимостью $c_j$, представляющей предельное число пользователей, для которых обслуживание ресурсом возможно без деградации качества. Каждый пользователь $i$ характеризуется профилем $p_i$ — описанием его интересов в форме, поддерживающей численное вычисление меры близости с описаниями ресурсов; в работе профиль представлен векторным эмбеддингом текстового описания интересов, согласованным по размерности с эмбеддингами описаний ресурсов.

Время рассматривается дискретно: эпизод взаимодействия пользователей с системой разбивается на тайм-слоты $T = \{1, \ldots, \tau\}$, в каждом из которых пользователю предъявляется рекомендательная подсказка и наблюдается фактический выбор. В пределах одного тайм-слота пользователь может быть сопоставлен не более чем с одним ресурсом; полный план взаимодействия пользователя с системой за эпизод формируется как последовательность сопоставлений по тайм-слотам.

\subsection{Управляющая политика}

Управляющая сторона реализует политику рекомендаций $\pi$, отображающую состояние системы в распределение возможных рекомендательных выдач. На каждом тайм-слоте $t$ для каждого пользователя $i$ политика формирует ранжированный список $L_{i,t}$ из $K$ ресурсов, отсортированных в порядке убывания предсказанной полезности. Список $L_{i,t}$ является непосредственным управляющим воздействием на пользователя; пользователь самостоятельно принимает решение о выборе ресурса, опираясь на представленный список и собственные предпочтения.

Существенным элементом постановки является зависимость рекомендации одного пользователя от рекомендаций, выдаваемых остальным пользователям в тот же тайм-слот. Формирование списка $L_{i,t}$ изолированно от $\{L_{j,t}\}_{j \neq i}$ воспроизводит классическую постановку индивидуальных рекомендательных систем и приводит, как показано в подразделе 1.2, к коллективным эффектам перегрузки. В рамках предлагаемой формализации политика $\pi$ принимает на вход агрегированное состояние системы, включающее текущую загрузку ресурсов, что обеспечивает информационную связность рекомендаций между пользователями.

\subsection{Релевантность как функция близости}

Полезность сопоставления пользователя $i$ с ресурсом $j$ определяется через меру семантической близости между профилем пользователя и описанием ресурса. Используется косинусная близость векторных эмбеддингов:

\[
r_{ij} = \cos(p_i, e_j) = \frac{p_i \cdot e_j}{\|p_i\| \cdot \|e_j\|},
\]

где $p_i$ — эмбеддинг профиля пользователя, $e_j$ — эмбеддинг описания ресурса. Выбор косинусной близости в качестве функции релевантности обусловлен её интерпретируемостью, низкой вычислительной стоимостью и широким распространением в качестве базового сигнала ранжирования в литературе рекомендательных систем.

Чисто семантическая близость в роли сигнала релевантности порождает известную проблему смыслового сужения выдачи (filter bubble): пользователь систематически получает рекомендации только тематически близких ресурсов. В работе эта проблема решается не на уровне определения релевантности, а на уровне сравниваемых политик: часть из них (в частности, MMR и LLM-ranker) предусматривают механизмы преодоления смыслового сужения, что позволяет провести систематическое сравнение в этом отношении.

\subsection{Отклик пользователя}

Ответ пользователя $i$ на рекомендательный список $L_{i,t}$ моделируется как случайная величина $y_{i,t}$, принимающая значения из множества $J \cup \{\varnothing\}$, где значение $\varnothing$ соответствует отказу от выбора в данном тайм-слоте. Распределение $y_{i,t}$ задаётся условной вероятностной функцией выбора $P(y_{i,t} \mid L_{i,t}, p_i, s_t)$, зависящей от рекомендательного списка $L_{i,t}$, профиля пользователя $p_i$ и текущего состояния системы $s_t$ (в первую очередь — текущей заполненности ресурсов). Вероятностный характер отклика отражает фундаментальное ограничение управляющего рычага: рекомендация формирует пользовательскую перцепцию доступных вариантов, но не предопределяет фактического выбора. Базовая функция выбора — мультиномиальная логит-модель с температурой: вероятность выбора ресурса пропорциональна экспоненте релевантности с учётом штрафа за заполненность, отказ представлен отдельной альтернативой с базовой вероятностью; точные параметры приводятся в главе 3.

Все компоненты совокупной награды рассчитываются по реализациям случайных величин $\{y_{i,t}\}$, а не по содержанию рекомендательных списков $\{L_{i,t}\}$: компонент пользовательской полезности агрегирует релевантности $r_{i, y_{i,t}}$ для тех пользователей и тайм-слотов, в которых $y_{i,t} \neq \varnothing$; мера переполнения и дисперсия заполненности рассчитываются по фактическому распределению пользователей по ресурсам, индуцированному реализациями $\{y_{i,t}\}$. Состояние $s_t$, наблюдаемое политикой на каждом шаге, аналогично формируется по реализациям $y_{i,t}$ для пользователей, обработанных ранее в текущем тайм-слоте. Тем самым проигнорированная пользователем рекомендация не приносит положительного вклада в полезность и не уменьшает наблюдаемой загрузки альтернативных ресурсов, что естественным образом штрафует политику за выдачу рекомендаций, не согласованных с предпочтениями пользователя.

\subsection{Целевая функция управляющей стороны}

Целью управляющей стороны является максимизация совокупного результата эпизода, агрегирующего три компонента: суммарную пользовательскую полезность выбранных ресурсов, минимизацию переполнений вместимости, минимизацию неравномерности загрузки ресурсов. Совокупная награда эпизода формируется как взвешенная сумма указанных компонентов:

\[
R = \alpha \cdot R_{\text{rel}} - \beta \cdot R_{\text{over}} - \gamma \cdot R_{\text{var}},
\]

где $R_{\text{rel}}$ — суммарная релевантность фактически выбранных пользователями ресурсов; $R_{\text{over}}$ — мера переполнения, агрегированная по тайм-слотам как сумма превышений вместимостей в долях от соответствующих лимитов; $R_{\text{var}}$ — дисперсия заполненности ресурсов, агрегированная по тайм-слотам; $\alpha, \beta, \gamma$ — неотрицательные весовые коэффициенты, определяющие относительную приоритетность компонентов. Конкретные значения коэффициентов определяются на этапе настройки и приводятся в главе 4.

\subsection{Постановка как марковский процесс принятия решений с ограничениями}

Описанная задача формализуется как марковский процесс принятия решений с ограничениями (Constrained MDP) в смысле, рассмотренном в подразделе 1.3.4. На каждом шаге $t$ управляющий агент наблюдает состояние $s_t$, включающее текущего обслуживаемого пользователя, его профиль, агрегированные характеристики уже обработанных пользователей и текущую загрузку ресурсов; выбирает действие $a_t = L_{i,t}$ — ранжированный список из $K$ ресурсов; получает мгновенную награду $r_t$, определяемую вкладом тайм-слота в совокупную награду $R$ согласно приведённой выше формуле; и переходит в новое состояние $s_{t+1}$, отражающее последствия выбора пользователя. Целью является нахождение стационарной политики $\pi^*$, максимизирующей ожидание совокупной награды по эпизоду.

\subsection{Связь с играми с перегрузкой}

Структурная природа задачи соответствует концепции игр с перегрузкой (congestion games) в смысле подраздела 1.3.2. Пользователи представляют автономных агентов, выбирающих ресурсы из предложенных вариантов; полезность ресурса для конкретного пользователя содержит компонент, зависящий от числа других пользователей, выбравших тот же ресурс (через ограничение вместимости и связанные с ним эффекты). В отсутствие управляющего вмешательства система движется к равновесию эгоистичных агентов, отдалённому от системного оптимума на величину, измеряемую показателем «цена анархии». Управляющая политика $\pi$ выступает в роли информационного сигнала, формирующего издержки альтернативных выборов в восприятии каждого пользователя и тем самым смещающего равновесие в направлении системно-оптимального распределения.

Предлагаемая постановка, таким образом, объединяет аппарат двух рамок: постановка задачи как Constrained MDP обеспечивает операциональный механизм построения и обучения управляющей политики, а трактовка через игру с перегрузкой — концептуальное обоснование необходимости совместной оптимизации рекомендаций между пользователями и теоретическую интерпретацию роли управляющей политики как механизма коррекции децентрализованного равновесия.

\section{Требования к решению}

Постановка задачи, сформулированная в подразделе 2.1, и анализ ограничений существующих подходов, выполненный в подразделе 1.5, позволяют сформулировать требования, которым должно удовлетворять разрабатываемое решение. Требования разделяются на три группы: функциональные требования, описывающие необходимые возможности системы; нефункциональные требования, описывающие её качественные характеристики; ограничения предметной области, определяющие условия, в которых решение должно работать.

\subsection{Функциональные требования}

К числу функциональных требований к разрабатываемому решению относятся:

\begin{enumerate}
  \item \textbf{Приём входных данных.} Система должна принимать на вход множество ресурсов с описаниями и вместимостями, а также множество пользователей с профилями интересов. Представление профилей и описаний должно поддерживать численное вычисление меры близости.
\end{enumerate}

\begin{enumerate}
  \item \textbf{Формирование персональных рекомендаций.} Для каждого пользователя в каждом тайм-слоте система должна формировать ранжированный список из $K$ ресурсов, упорядоченных по предсказанной полезности с учётом текущего состояния системы.
\end{enumerate}

\begin{enumerate}
  \item \textbf{Учёт ограничений вместимости.} Формирование рекомендаций должно учитывать вместимости ресурсов и текущую их заполненность; политика рекомендаций должна избегать систематического направления пользовательского потока к ресурсам, для которых превышение вместимости становится вероятным.
\end{enumerate}

\begin{enumerate}
  \item \textbf{Согласованность рекомендаций между пользователями.} Рекомендация, выдаваемая одному пользователю, должна формироваться с учётом рекомендаций, выданных и реализованных другими пользователями в текущем тайм-слоте.
\end{enumerate}

\begin{enumerate}
  \item \textbf{Обработка вероятностного отклика.} Система должна корректно обрабатывать стохастический характер пользовательского отклика, включая возможность выбора ресурса вне рекомендательного списка и отказа от выбора в тайм-слоте.
\end{enumerate}

\begin{enumerate}
  \item \textbf{Обновление состояния по фактическим выборам.} Состояние системы, наблюдаемое политикой, должно обновляться по реализованным пользовательским выборам, а не по выданным рекомендациям.
\end{enumerate}

\begin{enumerate}
  \item \textbf{Поддержка сравнения нескольких политик.} Система должна поддерживать прозрачное сравнение различных рекомендательных политик на одной и той же среде с использованием единого набора метрик.
\end{enumerate}

\subsection{Нефункциональные требования}

К числу нефункциональных требований относятся:

\begin{enumerate}
  \item \textbf{Воспроизводимость.} Прогон системы при фиксированных входных данных и параметрах среды должен давать одинаковый результат, что обеспечивается фиксацией начальных значений генераторов случайных чисел и явной спецификацией всех параметров эксперимента.
\end{enumerate}

\begin{enumerate}
  \item \textbf{Эффективность обучения.} Обучение и сравнение политик должны быть осуществимы в разумное время на доступном исследовательском оборудовании уровня одной графической вычислительной карты потребительского класса.
\end{enumerate}

\begin{enumerate}
  \item \textbf{Расширяемость.} Архитектура решения должна поддерживать добавление новых рекомендательных политик и новых функций отклика пользователей без переработки симуляционной среды и компонентов оценки качества.
\end{enumerate}

\begin{enumerate}
  \item \textbf{Прозрачность.} Все компоненты совокупной награды и итоговые метрики должны допускать содержательную интерпретацию в терминах задачи, что необходимо как для отладки, так и для последующего обоснования результатов работы.
\end{enumerate}

\begin{enumerate}
  \item \textbf{Низкая стоимость обучения.} Затраты на обучение и сравнение политик не должны существенно зависеть от стоимости обращения к большим языковым моделям. Использование больших языковых моделей допускается на этапах, требующих однократных или редких обращений, но не на каждом шаге обучения.
\end{enumerate}

\begin{enumerate}
  \item \textbf{Внешняя валидность результатов.} Качество результирующих политик должно подтверждаться не только в среде, в которой проходило обучение, но и в независимой среде с альтернативной моделью пользовательского поведения, что обеспечивает контроль разрыва между симуляцией и реальным поведением пользователей.
\end{enumerate}

\subsection{Ограничения предметной области}

Дополнительно фиксируются ограничения, определяющие условия, в которых решение должно сохранять работоспособность:

\begin{enumerate}
  \item \textbf{Малый каталог ресурсов.} Число рассматриваемых ресурсов исчисляется десятками, что исключает применение методов, требующих больших каталогов для устойчивого обучения базовых моделей релевантности.
\end{enumerate}

\begin{enumerate}
  \item \textbf{Отсутствие журналов реальных пользовательских взаимодействий.} Решение должно обеспечивать обучение и сравнение политик в условиях отсутствия исторических данных о фактических пользовательских выборах.
\end{enumerate}

\begin{enumerate}
  \item \textbf{Ограниченный бюджет вычислений и обращений к большим языковым моделям.} Объём обучающих и оценочных вычислений ограничен ресурсами одного исследователя; объём обращений к большим языковым моделям — соответствующим бюджетом.
\end{enumerate}

Совокупность сформулированных требований формирует основу для описания предлагаемого метода в подразделе 2.3.

\section{Описание предлагаемого метода}

Предлагаемый метод объединяет три компонента: гибридную симуляционную среду для обучения и оценки политик при отсутствии исторических данных; набор сравниваемых рекомендательных политик различной природы — от случайной до обучаемой методами обучения с подкреплением и основанной на большой языковой модели; двухслойную схему валидации с операциональной проверкой в обучающей среде и независимой проверкой в среде с альтернативной моделью пользовательского поведения. Каждый компонент адресует подмножество требований из подраздела 2.2.

\subsection{Гибридная симуляционная среда}

Симуляционная среда, в которой осуществляется обучение и сравнение рекомендательных политик, имеет двухкомпонентную структуру.

\textbf{Компонент однократной генерации пользовательской популяции.} Перед началом обучения большая языковая модель формирует синтетическую популяцию пользователей в количестве, достаточном для статистически значимого сравнения политик. Каждый пользователь характеризуется текстовым профилем интересов, отражающим его предметные предпочтения, уровень квалификации, тематическую специфику и иные характеристики, влияющие на пользовательский выбор. Сгенерированные профили преобразуются в векторные эмбеддинги, согласованные по размерности с эмбеддингами описаний ресурсов. Однократный характер обращения к большой языковой модели существенно снижает совокупную стоимость экспериментов по сравнению со схемами, в которых обращение к большой языковой модели выполняется на каждом шаге обучения, при сохранении содержательного разнообразия пользовательской популяции. Использование большой языковой модели в роли генератора синтетических данных описано в подразделе 1.4.4 как одно из основных направлений применения языковых моделей в задачах рассматриваемого класса.

\textbf{Компонент параметрического пошагового симулятора.} Основной цикл обучения и оценки политик реализуется параметрическим симулятором, разыгрывающим на каждом шаге функцию выбора пользователей. Функция выбора реализуется мультиномиальной логит-моделью с температурой, описанной в подразделе 2.1, и зависит от релевантности предложенных ресурсов профилю пользователя, текущей заполненности ресурсов и базовой склонности пользователя к отказу от участия в тайм-слоте. Параметры функции выбора (температура, штраф за заполненность, базовая вероятность отказа) фиксируются на этапе настройки симулятора. Параметрический характер компонента обеспечивает высокую скорость симуляции, воспроизводимость результатов при фиксации начальных значений генератора случайных чисел и независимость стоимости обучения от стоимости обращений к большой языковой модели.

Гибридная структура среды отвечает требованию низкой стоимости обучения (нефункциональное требование 5) при сохранении содержательного разнообразия популяции, недостижимого в чисто параметрических симуляторах. Кроме того, она удовлетворяет требованиям воспроизводимости и эффективности (нефункциональные требования 1 и 2).

\subsection{Множество сравниваемых рекомендательных политик}

В рамках предлагаемого метода рассматривается шесть рекомендательных политик различной природы, обеспечивающих систематическое сравнение различных подходов к решению задачи. Каждая политика реализует отображение текущего состояния системы и профиля обслуживаемого пользователя в ранжированный список из $K$ ресурсов.

\textbf{Политика 1: Случайный выбор.} Формирует ранжированный список путём равномерно случайного выбора $K$ ресурсов из каталога без учёта профиля пользователя и состояния системы. Не предполагает обучения. Назначение в составе сравнения — установление нижней планки качества, относительно которой содержательность остальных политик должна быть очевидна; при отсутствии существенного превышения над случайной политикой возникают основания усомниться в адекватности постановки задачи или метрик.

\textbf{Политика 2: Косинусная близость без учёта ограничений.} Формирует список как первые $K$ ресурсов, упорядоченных по убыванию косинусной близости между эмбеддингом профиля пользователя и эмбеддингами описаний ресурсов. Соответствует классическому подходу контентной фильтрации в рекомендательных системах без учёта системных ограничений. Назначение — представление точки отсчёта для подходов, не учитывающих ограничения вместимости и взаимосвязанность рекомендаций между пользователями.

\textbf{Политика 3: Косинусная близость с учётом вместимости (rule-based).} Расширение второй политики: к косинусной релевантности применяется мультипликативный штраф, монотонно убывающий с ростом текущей заполненности соответствующего ресурса; ресурсы, заполненные выше заданного порога, исключаются из ранжированного списка. Не предполагает обучения и реализует управление через детерминированное правило. Назначение — представление эвристического подхода, учитывающего ограничения вместимости без обращения к обучаемым методам.

\textbf{Политика 4: Maximal Marginal Relevance.} Реализует алгоритм MMR, описанный в подразделе 1.4.1, с балансом релевантности и разнообразия выдачи через параметр $\lambda$. Не предполагает обучения. Назначение — представление классического академического подхода к преодолению проблемы смыслового сужения выдачи; обеспечивает контроль того, насколько добавление разнообразия влияет на загрузку различных ресурсов в условиях рассматриваемой задачи.

\textbf{Политика 5: Обучаемая политика на основе обучения с подкреплением.} Реализуется методом проксимальной оптимизации политики (PPO), описанным в подразделе 1.4.2. Обучается в параметрическом симуляторе на совокупной награде, определённой в подразделе 2.1. Состояние, наблюдаемое политикой, включает профиль текущего обслуживаемого пользователя, текущую заполненность ресурсов и агрегированные характеристики уже обработанных пользователей. Действие политики — ранжированный список из $K$ ресурсов. Назначение — представление обучаемого подхода, способного учитывать структуру ограничений и согласованность рекомендаций между пользователями за счёт обучения непосредственно в симуляторе. Эта политика — основной технический результат работы.

\textbf{Политика 6: LLM-ranker.} Реализует подход, в котором большая языковая модель выступает в роли ранжирующего компонента. Модель получает на вход профиль пользователя, описания ресурсов из числа $N$ ресурсов с наибольшей косинусной близостью к профилю и сводку текущей заполненности залов; на выходе формирует ранжированный список из $K$ ресурсов с учётом всех представленных факторов. Соответствует подходу, описанному в подразделе 1.4.4 как одно из основных направлений применения больших языковых моделей в задачах рассматриваемого класса. Назначение — сравнение с современным альтернативным методом, представляющим парадигму непосредственного использования больших языковых моделей в качестве рекомендательного компонента.

Совокупность шести политик покрывает основные классы подходов, рассмотренные в подразделах 1.3 и 1.4: случайный выбор как формальная нижняя граница; чистое контентное ранжирование как классический recsys-подход без ограничений; правиловое управление с учётом ограничений; классический алгоритм диверсификации; обучаемая политика в постановке Constrained MDP; современный подход на основе большой языковой модели. Подобный охват обеспечивает систематическое сравнение и позволяет позиционировать обучаемую политику относительно как простых эвристик, так и современных альтернатив.

\subsection{Двухслойная схема валидации}

Внешняя валидность результатов обеспечивается двухслойной схемой проверки, разделяющей этапы обучения и сравнения политик и этап независимой проверки результата.

\textbf{Слой 1: основное сравнение в параметрическом симуляторе.} Все шесть политик обучаются (для политик 5 и 6 — настраиваются) и сравниваются в гибридной симуляционной среде, описанной выше. На этом слое реализуется протокол сравнения по совокупности метрик, определяемых в подразделе 2.4, и выполняется выбор лучшей политики по совокупному критерию. Среда первого слоя характеризуется высокой скоростью, воспроизводимостью и низкой стоимостью.

\textbf{Слой 2: независимая валидация лучшей политики в LLM-агентной среде.} Политика, признанная лучшей по результатам первого слоя, проходит дополнительную проверку в среде, в которой пользователи моделируются не параметрической мультиномиальной логит-моделью, а автономными агентами на основе больших языковых моделей с собственной памятью и функцией выбора, как описано в подразделе 1.4.4 со ссылкой на работу Agent4Rec. Среда второго слоя характеризуется существенно большей вычислительной стоимостью, что ограничивает её применение этапом итоговой проверки. Сопоставление результатов первого и второго слоёв обеспечивает оценку устойчивости выбранной политики к изменению модели пользовательского поведения, тем самым контролируя разрыв между симуляцией и реальным поведением (sim-to-real gap).

Двухслойная схема непосредственно адресует требование внешней валидности результатов: разделение слоёв по стоимости обеспечивает совместимость требований внешней валидности и низкой стоимости обучения, которые в одноуровневой схеме противоречат друг другу.

\subsection{Соответствие требованиям}

Предлагаемый метод покрывает все требования, сформулированные в подразделе 2.2. Функциональные требования 1–7 обеспечиваются совокупностью симуляционной среды (требования 1, 5, 6) и сравниваемых политик (требования 2, 3, 4, 7), причём учёт ограничений вместимости и согласованность рекомендаций между пользователями реализуются на уровне состояния, наблюдаемого политикой. Нефункциональные требования удовлетворяются гибридной структурой среды (требования 1, 2, 5), архитектурой раздельно обучаемых политик (требование 3), интерпретируемой структурой совокупной награды (требование 4) и двухслойной схемой валидации (требование 6). Ограничения предметной области учитываются за счёт переноса основной нагрузки обучения в параметрический симулятор (ограничение 3), синтетической генерации пользовательской популяции (ограничение 2) и поддержки малого каталога ресурсов всеми компонентами метода (ограничение 1).

\section{Сравнительный анализ альтернатив}

Помимо предлагаемого метода, описанного в подразделе 2.3, существуют альтернативные методологические подходы, опирающиеся на другие концептуальные рамки из числа рассмотренных в подразделе 1.3 либо на иные конфигурации симуляционных сред. Настоящий подраздел проводит сравнительный анализ четырёх альтернатив с предлагаемым методом и обосновывает выбор последнего.

\subsection{Альтернатива 1. Чисто оптимизационный подход в рамках исследования операций}

Альтернативой предлагаемому методу является формулировка задачи как обобщённой задачи о назначениях (GAP) или многопериодной задачи о назначениях, рассмотренной в подразделе 1.3.1, с последующим её решением точными методами целочисленного программирования или приближёнными эвристиками. Подобный подход обеспечивает строгую математическую постановку и развитый алгоритмический арсенал.

Существенным препятствием для применения данной альтернативы в качестве основного метода является фундаментальное допущение, лежащее в основе классической постановки задач о назначениях: решение управляющей стороны полагается исполнимым, то есть приравнивается к фактическому распределению пользователей по ресурсам. В рассматриваемой задаче решение управляющей стороны имеет статус рекомендации, которой пользователь может следовать или не следовать. Прямое применение подхода исследования операций к задаче рекомендательного управления приводит к решениям, оптимальным в детерминированной модели, но потенциально неоптимальным с учётом вероятностного отклика пользователей.

Вместе с тем, оптимизационный подход сохраняет ценность в составе предлагаемого метода: решение детерминированной задачи о назначениях используется как теоретическая верхняя граница достижимой суммарной полезности при идеальном следовании пользователей рекомендациям, что обеспечивает контрольную точку для интерпретации результатов сравниваемых политик.

\subsection{Альтернатива 2. Чисто теоретико-игровой подход}

Альтернативой является описание задачи как игры с перегрузкой в рамках теоретико-игровой рамки подраздела 1.3.2 с последующим аналитическим вычислением равновесия Нэша или равновесия Уордропа, а также показателя «цена анархии» для системы. Подобный подход опирается на развитую теорию и обеспечивает концептуально чистое описание ситуации децентрализованного выбора автономных агентов.

Препятствием для применения данной альтернативы в качестве основного метода является описательный, а не предписывающий характер классической теоретико-игровой рамки. Аппарат равновесий отвечает на вопрос, какое распределение установится при заданных правилах взаимодействия и функциях издержек, однако не предоставляет алгоритма построения конкретного управляющего воздействия — рекомендательной выдачи, — способного сместить поведение пользователей в направлении системного оптимума. Переход от описания равновесий к проектированию управляющих воздействий требует обращения к смежным теоретическим направлениям, развитым преимущественно для целей теоретического анализа и не предоставляющим готовых вычислительных алгоритмов.

В предлагаемом методе теоретико-игровая рамка используется как структурное обоснование необходимости совместной оптимизации рекомендаций между пользователями (подраздел 2.1) и теоретическая интерпретация роли управляющей политики, но операциональным методом построения политики выступает аппарат марковского процесса принятия решений с ограничениями.

\subsection{Альтернатива 3. Прямое обучение в LLM-агентной среде}

Альтернативой является использование LLM-агентной симуляционной среды, описанной в подразделе 1.4.4, в качестве основной среды для обучения рекомендательной политики, без параметрического симуляционного слоя. Подобная среда обеспечивает максимальную содержательную глубину пользовательского поведения, моделируя каждого пользователя как автономного агента на основе большой языковой модели с собственной памятью и функцией выбора.

Препятствием для применения данной альтернативы является вычислительная стоимость. Обучение рекомендательной политики методами обучения с подкреплением требует значительного числа взаимодействий со средой — порядка нескольких миллионов шагов в типовых конфигурациях. В LLM-агентной среде каждый шаг включает множественные обращения к большой языковой модели, что превращает совокупный бюджет обучающего эксперимента в непрактично высокий при ограниченных вычислительных ресурсах.

В предлагаемом методе LLM-агентная среда применяется как слой 2 валидации, охватывающий только итоговую проверку лучшей политики на ограниченном числе сценариев, что укладывается в реалистичный экспериментальный бюджет.

\subsection{Альтернатива 4. Прямое применение классических методов рекомендательных систем с ограничениями}

Альтернативой является прямое применение методов рекомендательных систем с ограничениями, описанных в подразделе 1.3.3, — в первую очередь, постановки ранжирования с ограничениями Сингха и Йоахимса. Указанные методы обеспечивают формальный аппарат и алгоритмы оптимизации политики ранжирования при дополнительных ограничениях на агрегированную экспозицию ресурсов.

Препятствием для применения данной альтернативы является принципиальное отличие в типе формулируемого ограничения. Классические методы рекомендательных систем с ограничениями формулируют ограничения как математическое ожидание по распределению пользователей, что подразумевает выравнивание совокупной экспозиции ресурсов «в среднем» по большому потоку независимых выдач. В рассматриваемой постановке ограничения вместимости должны соблюдаться в каждой конкретной коллективной выдаче, формируемой одновременно для фиксированного множества пользователей. Адаптация классических методов под подобный режим работы требует существенной методологической переработки, по объёму сопоставимой с разработкой нового метода.

Вместе с тем, методы данного направления находят применение в составе сравниваемых политик предлагаемого метода — в частности, политика 4 (Maximal Marginal Relevance) реализует один из канонических алгоритмов переранжирования, описанных в подразделе 1.4.1.

\subsection{Сводное обоснование выбора}

Сравнительный анализ показывает, что ни одна из рассмотренных альтернатив, взятая отдельно, не закрывает совокупности требований, сформулированных в подразделе 2.2. Подход исследования операций (альтернатива 1) не учитывает вероятностного отклика; теоретико-игровой подход (альтернатива 2) описателен, а не предписывающ; прямое обучение в LLM-агентной среде (альтернатива 3) непрактично по стоимости; прямое применение методов рекомендательных систем с ограничениями (альтернатива 4) опирается на формализм амортизированных, а не одноразовых ограничений.

Предлагаемый метод — синтез элементов нескольких рамок: операциональная постановка взята из аппарата марковского процесса принятия решений с ограничениями (подраздел 1.3.4), концептуальное обоснование совместной оптимизации опирается на теорию игр с перегрузкой (подраздел 1.3.2), сравниваемые политики покрывают разные классы подходов из подразделов 1.3.3, 1.4.1 и 1.4.4, двухслойная схема валидации согласует требования низкой стоимости обучения и внешней валидности результатов. Такая конфигурация удовлетворяет всем требованиям одновременно, что не достижимо ни одной из альтернатив в её прямой форме.

\section{Дизайн экспериментов и план валидации}

Настоящий подраздел определяет дизайн экспериментальной части работы, направленной на сравнительную оценку шести политик, описанных в подразделе 2.3, и независимую валидацию лучшей политики во втором слое симуляционной среды. Экспериментальная часть включает четыре составляющие: сравнение политик по основным метрикам качества, исследования влияния отдельных компонентов метода (ablation studies), проверку устойчивости результатов к вариации параметров среды и итоговую валидацию в LLM-агентной среде. Все экспериментальные процедуры опираются на единый набор метрик, описанный в первой части подраздела.

\subsection{Основные метрики оценки}

Сравнение рекомендательных политик опирается на четыре метрики, каждая из которых отражает отдельный аспект качества решения.

\textbf{Overflow rate.} Относительная мера переполнений ресурсов. Для каждого ресурса в каждом тайм-слоте вычисляется превышение фактического числа пользователей над вместимостью, отнесённое к вместимости; итоговое значение — среднее этих превышений по всем ресурсам, тайм-слотам и эпизодам. Метрика принимает значения от нуля (вместимости соблюдаются строго) до положительных значений, растущих с увеличением степени переполнения. Отражает соблюдение жёстких ограничений вместимости в каждой одноразовой выдаче.

\textbf{Hall utilization variance.} Дисперсия относительной заполненности ресурсов в пределах тайм-слота, усреднённая по тайм-слотам и эпизодам. Низкие значения метрики соответствуют равномерному распределению аудитории между ресурсами; высокие значения — ситуации одновременной перегрузки части ресурсов и недозагрузки других. Отражает согласованность рекомендаций, выдаваемых разным пользователям одновременно.

\textbf{Mean user utility.} Средняя релевантность ресурсов, фактически выбранных пользователями в ходе эпизода. Релевантность измеряется как косинусная близость профиля пользователя и описания выбранного ресурса (подраздел 2.1); пользователи, отказавшиеся от выбора в тайм-слоте, вносят нулевой вклад в среднее. Метрика отражает суммарную пользовательскую полезность, обеспечиваемую политикой.

\textbf{Gini coefficient по ресурсам.} Коэффициент Джини, рассчитанный по совокупному распределению аудитории между ресурсами за эпизод, отражающий неравенство получаемого ресурсами внимания. Близкие к нулю значения соответствуют равномерному распределению аудитории; близкие к единице — ситуации, в которой большая часть пользователей сосредоточена на малом числе ресурсов. Метрика обеспечивает оценку справедливости со стороны поставщиков ресурсов в смысле, рассмотренном в подразделе 1.3.3.

Совокупность четырёх метрик покрывает три аспекта качества: соблюдение системных ограничений (Overflow rate, Hall utilization variance), пользовательская полезность (Mean user utility), справедливость распределения внимания между ресурсами (Gini coefficient). Раздельное представление метрик обеспечивает возможность анализа компромиссов, возникающих между конкурирующими аспектами качества.

\subsection{Протокол сравнения политик}

Сравнение шести политик из подраздела 2.3 проводится в первом слое симуляционной среды по следующему протоколу.

В качестве экспериментальной обстановки фиксируется конкретная конфигурация среды: множество ресурсов с заданными описаниями и вместимостями, синтетическая популяция пользователей, сформированная однократной LLM-генерацией, расписание тайм-слотов. Все шесть политик прогоняются в идентичной среде; различие между прогонами определяется только начальным значением генератора случайных чисел, управляющего стохастикой функции выбора.

Каждая политика оценивается на множестве независимых эпизодов количества $n_{\text{ep}}$, отличающихся начальным значением случайных чисел. Для обучаемых политик (политика 5) обучение выполняется отдельно с использованием фиксированной части эпизодов, оценка — на отдельной отложенной выборке, что обеспечивает корректное разделение обучающих и оценочных данных. Для каждой метрики вычисляются среднее и доверительный интервал по эпизодам.

Сравнение результатов проводится по двум критериям. Первый критерий — раздельное сравнение по каждой из четырёх метрик; политики ранжируются по каждой метрике с указанием статистической значимости различий. Второй критерий — совокупное сравнение по взвешенной сумме компонентов совокупной награды, определённой в подразделе 2.1; этот критерий используется для выбора лучшей политики при наличии компромиссов между метриками.

\subsection{Исследование влияния отдельных компонентов метода (ablation studies)}

С целью оценки относительного вклада отдельных элементов предлагаемого метода в итоговое качество проводится серия исследований, в каждом из которых отключается или заменяется один из компонентов, после чего повторно проводится сравнение политик.

\textbf{Влияние LLM-генерации синтетической популяции.} Проводится сравнение результатов основного эксперимента с результатами эксперимента, в котором профили пользователей формируются равномерно случайно из пространства эмбеддингов либо отбираются случайно из заданного фиксированного набора. Цель — оценить, в какой мере содержательное разнообразие популяции, обеспечиваемое большой языковой моделью, сказывается на итоговых результатах сравнения политик.

\textbf{Влияние компонентов совокупной награды.} Проводится повторное обучение политики 5 (обучаемой методом проксимальной оптимизации политики) в трёх режимах: с полной совокупной наградой; с отключённым штрафом за переполнение ($\beta = 0$); с отключённым штрафом за дисперсию заполненности ($\gamma = 0$). Цель — определить относительный вклад каждой компоненты награды в результирующее поведение политики.

\textbf{Влияние информационной связности состояния.} Проводится повторное обучение политики 5 в режиме, в котором состояние, наблюдаемое политикой, не включает информацию о текущей заполненности ресурсов. Цель — эмпирически подтвердить значимость информационной связности рекомендаций между пользователями для решения задачи; ожидается, что отключение данной составляющей состояния приводит к существенной деградации метрик соблюдения вместимости.

\subsection{Проверка устойчивости результатов}

Устойчивость результатов сравнения политик к вариации параметров среды проверяется серией экспериментов, в которых варьируется один из параметров среды при фиксации остальных:

\begin{itemize}
  \item Объём пользовательской популяции — варьируется в диапазоне нескольких порядков для оценки масштабируемости методов;
  \item Распределение вместимостей ресурсов — рассматриваются равномерное и существенно неравномерное распределения;
  \item Параметр температуры функции выбора в симуляторе — варьируется в диапазоне, соответствующем переходу от детерминированного следования рекомендациям до близкого к равномерному выбору;
  \item Размер каталога ресурсов — варьируется в пределах, оставляющих каталог в режиме малого размера в смысле подраздела 2.2.
\end{itemize}

Для каждой комбинации параметров повторяется протокол сравнения политик. Согласованность ранжирования политик при различных значениях параметров рассматривается как индикатор устойчивости предлагаемого метода.

\subsection{Итоговая валидация во втором слое}

Лучшая политика по результатам первого слоя проходит независимую валидацию во втором слое симуляционной среды на основе LLM-агентов, описанной в подразделе 2.3. Протокол валидации включает прогон политики на ограниченном числе эпизодов в среде, в которой пользователи моделируются автономными агентами на основе больших языковых моделей с собственной памятью и функцией выбора.

В рамках второго слоя рассчитываются те же четыре метрики, что и в первом, и проводится их сопоставление с результатами первого слоя. Отдельное внимание уделяется качественному анализу расхождений: сценариям, в которых результаты двух слоёв существенно различаются. Подобный анализ обеспечивает оценку разрыва между симуляцией и реальным поведением (sim-to-real gap) и формирует материал для обсуждения внешней валидности результатов работы.

\subsection{Критерии успешности работы}

Совокупность экспериментов считается подтверждающей предлагаемый метод при выполнении трёх условий:

\begin{enumerate}
  \item Обучаемая политика (политика 5) демонстрирует статистически значимое превосходство над политиками 1–4 по совокупному критерию и по большинству отдельных метрик в первом слое симуляционной среды.
  \item Обучаемая политика 5 показывает результаты, сопоставимые с результатами политики 6 (LLM-ranker), либо превосходит её по совокупному критерию при существенно меньшей вычислительной стоимости отдельной рекомендации.
  \item Итоговая валидация во втором слое качественно подтверждает результаты первого слоя; наблюдаемые расхождения метрик находятся в пределах, не отменяющих основных выводов о ранжировании политик.
\end{enumerate}

Невыполнение одного или нескольких условий рассматривается как самостоятельный результат работы, подлежащий обсуждению в главе 4 в части анализа полученных результатов и формулировки направлений дальнейших исследований.

\section{Выводы по главе}

В главе сформулирована постановка задачи и описан предлагаемый методологический подход к её решению.

В подразделе 2.1 проведена формализация задачи на основе синтеза двух концептуальных рамок, рассмотренных в главе 1: марковского процесса принятия решений с ограничениями и игр с перегрузкой. Введены формальные обозначения объектов задачи, описана структура управляющей политики, определена функция релевантности на основе косинусной близости эмбеддингов профилей и описаний ресурсов, формализован вероятностный отклик пользователей через случайные величины фактических выборов. Целевая функция управляющей стороны представлена как взвешенная сумма пользовательской полезности, штрафа за переполнение вместимостей и штрафа за неравномерность загрузки ресурсов.

В подразделе 2.2 сформулированы требования, которым должно удовлетворять разрабатываемое решение. Требования разделены на три группы: семь функциональных требований к возможностям системы, шесть нефункциональных требований к её качественным характеристикам и три ограничения предметной области.

В подразделе 2.3 описан предлагаемый метод, объединяющий три компонента: гибридную симуляционную среду с однократной LLM-генерацией пользовательской популяции и параметрическим пошаговым симулятором; множество из шести сравниваемых рекомендательных политик различной природы; двухслойную схему валидации с основным сравнением в параметрическом симуляторе и независимой проверкой лучшей политики в среде на основе LLM-агентов. Показано, что предложенная конфигурация метода удовлетворяет всем требованиям, сформулированным в подразделе 2.2.

В подразделе 2.4 проведён сравнительный анализ четырёх альтернативных методологических подходов: чисто оптимизационного подхода в рамках исследования операций; чисто теоретико-игрового подхода; прямого обучения политики в LLM-агентной среде; прямого применения классических методов рекомендательных систем с ограничениями. Установлено, что ни одна из альтернатив, взятая отдельно, не закрывает совокупности требований, что обосновывает выбор предлагаемого метода как систематического синтеза элементов нескольких рамок.

В подразделе 2.5 определён дизайн экспериментальной части работы, включающий четыре составляющие: сравнение политик по основным метрикам качества (Overflow rate, Hall utilization variance, Mean user utility, Gini coefficient); исследования влияния отдельных компонентов метода (LLM-генерация персон, компоненты совокупной награды, информационная связность состояния); проверку устойчивости результатов к вариации параметров среды; итоговую валидацию лучшей политики во втором слое симуляционной среды. Сформулированы три критерия успешности работы.

На основании сформулированной постановки задачи, требований и описания метода в главе 3 представляются конкретные технические решения по реализации симуляционной среды и сравниваемых политик, а также описывается ход проведённых экспериментов и полученные результаты.

\chapter{РЕАЛИЗАЦИЯ ПРОГРАММНОГО КОМПЛЕКСА}
\label{ch:implementation}
\sloppy


\section{Архитектура реализованного программного комплекса}
\label{sec:3-1}

Глава~\ref{ch:method} определяет состав системы как набор из девяти функциональных модулей, образующих систему поддержки принятия решений для организатора программы конференции. Настоящий подраздел описывает реализацию этого замысла в виде программного комплекса.

Комплекс собран из нескольких модулей: модели, симуляторы, рекомендательные политики, оператор локальных модификаций программы, планировщик экспериментов, модуль показателей и средства формирования сравнительных отчётов. Модули взаимодействуют через стабильные интерфейсы; это даёт воспроизводимость отдельных этапов эксперимента, независимое тестирование модулей и расширение состава политик и моделей популярности без переработки симулятора.

В реализации комплекс разделён на три архитектурных уровня. Уровень данных хранит и подготавливает входы: программа конференции, синтетическая популяция профилей участников, векторные представления описаний докладов и профилей, преднасчитанная оценка априорной популярности докладов. Уровень моделирования реализует отклик аудитории на программу и на рекомендации: модель поведения участника, семейство политик рекомендаций, оператор локальных модификаций программы и два симулятора с различной формой модели выбора --- параметрический и агентский на основе языковой модели. Уровень эксперимента отвечает за дизайн численного эксперимента и обработку его результатов: генератор плана эксперимента, управление случайностью на принципе общих случайных чисел, модуль показателей загрузки и риска, средства формирования сравнительного отчёта. Зависимости между уровнями однонаправленны: данные не знают ничего о моделях и эксперименте; моделирование зависит от данных, но не от плана эксперимента; эксперимент собирает результаты, но не вмешивается в логику отклика. Программа конференции остаётся единственным фактическим входом системы; все остальные сущности либо моделируются, либо задаются параметрами эксперимента.

Разработанный программный комплекс рассматривается как аналитический слой существующей продуктовой экосистемы для конференций сообщества JUG, на стороне организатора. Клиентская часть экосистемы --- приложение для индивидуального формирования участником расписания докладов --- была развёрнута в производственной среде и обеспечивает существующий канал взаимодействия с участниками. Комплекс подключается к этой экосистеме как самостоятельный аналитический инструмент: на вход он принимает программу конференции в формате, совместимом с существующим сервисом рекомендаций, и возвращает сравнительный отчёт для организатора. Метрики использования клиентского приложения и пользовательские сигналы в постановку настоящей работы не входят и комплексом не используются.

Реализация программного комплекса выполнена автором лично. К авторским инженерным решениям относятся: разделение каналов влияния политики и состояния системы на функцию полезности участника; форма функции полезности с симплексным соотношением весов и реализованные представители её компонентов; форма оператора локальных модификаций программы и протокол его применения; паритет каналов между параметрическим и агентским симуляторами как принцип проектирования контекста для языковой модели; схема разделения источников случайности при общих случайных числах; протокол прохождения граничной верификации модели как блокирующего фильтра перед содержательными выводами. Существующий клиентский сегмент продуктовой экосистемы JUG используется как контекст интеграции и в состав разработанного комплекса не входит.

Три архитектурных решения определяют комплекс в целом и обсуждаются ниже. Первое --- разделение каналов влияния политики и состояния системы на функцию полезности участника: политика влияет на распределение посещений строго через компонент функции полезности с весом $w_{rec}$, а ограничения вместимости учитываются на уровне политики, не на уровне выбора участника. Благодаря этому требование о граничном поведении модели --- при $w_{rec} \to 0$ политики неразличимы --- выполняется по построению, а риск двойного учёта вместимости снят. Подробное обсуждение реализации модели поведения вынесено в подраздел~\ref{sec:3-3}.

Второе решение --- двусимуляторная архитектура моделирования отклика. Параметрический симулятор работает в замкнутой форме выбора и используется как основной инструмент полного прогона по плану эксперимента; агентский симулятор на основе языковой модели работает на репрезентативном подмножестве конфигураций и используется для перекрёстной проверки выводов о сравнительной устойчивости политик. Оба симулятора реализованы как взаимозаменяемые компоненты, работающие через единый интерфейс с модулем программы и модулем политик. Это удовлетворяет защищаемому положению о взаимной проверке моделей; реализация описана в подразделе~\ref{sec:3-5}.

Третье решение --- построение плана эксперимента на основе латинского гиперкуба с применением общих случайных чисел внутри одной точки выборки. План эксперимента отделён от логики симуляции и реализован как самостоятельный модуль, формирующий выборку конфигураций по заданным параметрическим осям. При общих случайных числах внутри одной точки гиперкуба все политики применяются к одной и той же синтетической аудитории и одной и той же реализации варианта программы; шум сравнения политик при этом существенно снижается. Реализация описана в подразделе~\ref{sec:3-6}.

Дальнейшие подразделы настоящей главы последовательно раскрывают сценарий выполнения комплекса (подраздел~\ref{sec:3-2}), реализацию модели отклика участника (подраздел~\ref{sec:3-3}), реализацию управляющих воздействий --- политик рекомендаций и оператора локальных модификаций программы (подраздел~\ref{sec:3-4}), агентский симулятор как второй источник отклика (подраздел~\ref{sec:3-5}), реализацию плана эксперимента и обработки результатов (подраздел~\ref{sec:3-6}), использованные данные и верификацию работоспособности (подраздел~\ref{sec:3-7}), технологический стек и порядок использования систем искусственного интеллекта (подраздел~\ref{sec:3-8}). Подраздел~\ref{sec:3-9} содержит выводы по главе.


\section{Пайплайн выполнения сценарного эксперимента}
\label{sec:3-2}

Сценарный эксперимент устроен как набор однотипных единичных итераций. Единичная итерация --- вычисление показателей загрузки, риска перегрузки и релевантности для одной комбинации <<программа конференции --- конфигурация модели --- политика рекомендаций --- реализация случайности>>. Подготовка данных, перебор конфигураций и агрегация результатов обслуживают этот единичный прогон.

На входе фиксируется программа конференции --- структурированное описание слотов, докладов, залов и их номинальных вместимостей, а также метаданные докладов: тексты аннотаций, тематические признаки, состав спикеров. Номинальная вместимость отражает физические свойства площадки и сама по себе в исходных данных не варьируется; её масштабирование в сценарных конфигурациях выполняется отдельной осью эксперимента, описанной ниже, и применяется в момент прогона без изменения исходной программы. Программа загружается из файла фиксированного формата; векторные представления описаний докладов и профилей участников считаются заранее, на этапе подготовки данных, и хранятся отдельно. Это разделение продиктовано двумя соображениями. Первое --- векторные представления получены однократным вызовом внешней модели представления текстов и не должны пересчитываться при каждом запуске эксперимента. Второе --- отделение представлений от программы делает возможным независимую замену источника эмбеддингов без изменения логики симулятора.

Параметры сценарного эксперимента формируются генератором плана эксперимента. План задаёт набор точек в пространстве параметрических осей конфигурации, перечисленных в главе~\ref{ch:method}: вместимость залов, источник оценки популярности докладов, веса модели поведения участника, размер аудитории, вариант программы. Внутри каждой точки гиперкуба перебираются все четыре политики рекомендаций; каждая комбинация <<точка $\times$ политика>> повторяется на нескольких независимых случайных зёрнах для усреднения. План построен на латинском гиперкубе, дающем равномерное покрытие пространства параметров заданным числом точек при существенном сокращении объёма по сравнению с полным факторным перебором. Подробное обсуждение этого решения вынесено в подраздел~\ref{sec:3-6}.

Перед запуском единичной итерации система формирует синтетическую аудиторию. Аудитория получается отбором из заранее подготовленной популяции профилей: это сокращает стоимость формирования и обеспечивает воспроизводимость. Размер аудитории задаётся осью эксперимента, состав отбора --- отдельным случайным зерном, фиксированным на уровне точки гиперкуба и одинаковым для всех политик внутри этой точки. Сравнение политик внутри одной точки выполняется на одной и той же синтетической аудитории. Это реализация принципа общих случайных чисел из методики планирования имитационных экспериментов; она снимает шум сравнения политик и делает попарные сопоставления внутри точки информативными.

Источник оценки популярности докладов выбирается параметром конфигурации. Реализованы три варианта: оценка только по содержательной близости описания доклада к тематике конференции, оценка только по структурным признакам известности доклада в программе и комбинированная оценка. Соответствие источника популярности и поведения участника задаётся через функцию полезности и фиксируется в подразделе~\ref{sec:3-3}.

Вариант программы --- это либо исходная программа, либо одна из её допустимых локальных модификаций, полученных оператором перестановок докладов между разными временными слотами одного дня. Реализация оператора и проверка ограничения о конфликтах спикеров описаны в подразделе~\ref{sec:3-4}. Множество допустимых вариантов программы фиксируется на уровне точки гиперкуба и одинаково для всех политик внутри этой точки; это продолжает принцип общих случайных чисел и делает сопоставимыми результаты сравнения политик при одной и той же реализации структуры программы.

Единичный прогон симулятора обрабатывает программу слот за слотом. Внутри одного временного слота участники обрабатываются последовательно: каждый следующий участник наблюдает обновлённое состояние загрузки залов, отражающее решения тех, кто был обработан раньше. Это даёт реалистичный эффект кумулятивного формирования распределения посещений. Шаг для одного участника состоит из трёх связанных действий: политика формирует упорядоченный список рекомендаций; модель поведения участника вычисляет распределение вероятностей выбора над докладами слота; симулятор разыгрывает выбор и обновляет состояние загрузки. Логика модели поведения и логика политики разделены явно: политика влияет на выбор участника только через компонент функции полезности с соответствующим весом, как зафиксировано в постановке.

После прохождения всех слотов программы для одной комбинации параметров система фиксирует протокол шагов и состояние загрузки, накопленное к концу каждого слота. На основе этих данных модуль показателей вычисляет метрики из набора, заданного в главе~\ref{ch:method}. Показатели сохраняются в табличном формате с явным указанием значения каждой оси конфигурации и идентификатора политики, что даёт возможность последующего попарного сравнения и анализа чувствительности.

Воспроизводимость прогонов держится на строгом разделении источников случайности. Случайность в системе возникает в трёх местах: при отборе синтетической аудитории, при формировании множества вариантов программы оператором локальных модификаций и при разыгрывании выбора участника в симуляторе. Каждый из этих источников управляется отдельным случайным зерном с явным правилом образования. Зерно отбора аудитории и зерно оператора модификаций программы фиксируются по идентификатору точки гиперкуба и одинаковы для всех политик и реплик внутри этой точки; зерно симулятора фиксируется по номеру реплики и варьируется между ними, что и даёт независимое усреднение. Внутри самого симулятора разделены поток случайности окончательного выбора участника и поток случайности стохастической части политики. Это делает граничное поведение модели строго воспроизводимым: при нулевом весе компонента рекомендации выбор участника не зависит от политики --- свойство, которое глава~\ref{ch:method} относит к обязательным для граничной верификации.

Завершение единичной итерации даёт строку в табличном формате результатов. Полный сценарный эксперимент состоит из последовательного выполнения всех таких итераций по плану эксперимента и их объединения в единую таблицу. Дальнейшая обработка результатов --- попарные сравнения политик, оценка устойчивости при варьировании одной из осей конфигурации, эффект локальных модификаций программы относительно исходной программы и эффект социального заражения в зависимости от параметров --- выполняется отдельным модулем постобработки на уже сформированной таблице. Результаты этой обработки сохраняются в виде отдельных артефактов и подаются на вход модулю формирования сравнительного отчёта. Численные результаты постобработки и их интерпретация представлены в главе~\ref{ch:results}.


\section{Реализация модели отклика участника}
\label{sec:3-3}

Модель отклика участника --- центральный компонент параметрического симулятора. На каждом шаге симуляции она получает на вход профиль участника, множество доступных в текущем временном слоте докладов, упорядоченный список рекомендаций политики и текущее состояние загрузки залов; на выходе задаёт распределение вероятностей выбора над докладами слота, по которому симулятор разыгрывает фактический выбор. Все сравнительные результаты эксперимента определяются формой этой модели и значениями её параметров.

Модель реализована в форме мультиномиальной логит-модели с аддитивной функцией полезности из трёх компонентов --- содержательной релевантности, следования рекомендации и социального заражения; стохастичность выбора задаётся отдельным параметром температуры. Эти компоненты введены в постановке главы~\ref{ch:method}; здесь рассматривается их инженерная реализация. В принятой реализации в качестве канала следования рекомендации использована индикаторная форма, а в качестве компонента социального заражения --- логарифмическая нормировка по числу выбравших доклад участников; оба представителя выбраны как простейшие содержательные формы соответствующих классов и подробно обсуждаются ниже. Веса трёх компонентов связаны симплексным ограничением: их сумма равна единице. Симплексное ограничение реализовано на уровне генератора плана эксперимента, не на уровне модели поведения; модель поведения принимает три веса как независимые параметры. Это упрощает интерпретацию весов как долей влияния трёх каналов и сокращает размерность пространства параметрических осей, не теряя содержания постановки.

Реализация опирается на четыре инженерных решения, каждое из которых обсуждается ниже.

\paragraph{Учёт вместимости как ответственность политики, не модели поведения.}
В первой реализации параметрического симулятора, выполненной до фиксации текущей постановки, ограничение вместимости было зашито в функцию полезности участника как штрафной член, монотонно убывающий по текущей загрузке зала. Это решение пересмотрено по двум содержательным причинам.

Первая причина --- противоречие с принципом разделения каналов, принятым в главе~\ref{ch:method}: политика рекомендаций влияет на распределение посещений строго через компонент функции полезности с весом $w_{rec}$. Если ограничение вместимости входит в функцию полезности напрямую, оно действует в обход этого канала: даже контрольная политика, не предъявляющая участнику никакой рекомендации, оказывается косвенно подвержена эффекту вместимости. Это нарушает разделение каналов и делает невозможным выполнение свойства неразличимости политик при $w_{rec} \to 0$, обязательного по требованию граничной верификации модели.

Вторая причина --- расхождение с практикой, наблюдаемой в литературе. В рассмотренных в главе~\ref{ch:domain} работах по рекомендательным системам с учётом ограничений вместимости вместимость трактуется как свойство стороны рекомендатора, а не стороны участника. Метод ReCon реализует учёт через целевую функцию обучения рекомендатора с компонентом транспортной стоимости; формализм FEIR --- через переранжирование выдачи; capacity-aware методы в области рекомендации точек интереса --- через постпроцессинговый шаг распределения. Во всех трёх случаях участник остаётся моделью индивидуального выбора без явного знания о вместимости.

В принятой реализации ограничение вместимости полностью изъято из функции полезности участника. Учёт вместимости целиком ложится на одну из политик семейства --- политику с учётом загрузки залов, описанную в подразделе~\ref{sec:3-4}. Эта политика добавляет к скорингу кандидатов компонент, монотонно убывающий по текущей загрузке зала, и отбрасывает залы с загрузкой выше заданного жёсткого порога. Так достигается симметричное соответствие постановке: политика влияет на выбор участника только через компонент полезности с весом $w_{rec}$, а ограничение вместимости становится одним из возможных дизайнов политики и сравнивается с альтернативами на общих основаниях.

\paragraph{Канал рекомендации как индикатор присутствия в выдаче.}
Канал следования рекомендации в принятой реализации задан индикаторной функцией: компонент полезности равен единице, если доклад входит в упорядоченный список рекомендаций, выданный политикой, и нулю в противном случае. Это самая простая из возможных форм, выбрана сознательно.

Альтернативные формы --- ранговая, в которой вклад убывает с позицией доклада в списке, и нормированная по скорингу политики --- рассматривались на этапе проектирования модели. Содержательно ни одна из них не была отвергнута, но в условиях отсутствия данных о фактическом поведении аудитории выбор более сложной формы означает введение в модель дополнительной свободы без возможности её калибровки. Индикаторная форма зафиксирована как минимальная содержательная: она удовлетворяет требованию о граничном поведении модели по $w_{rec}$ и оставляет возможность последующего расширения без переписывания общей структуры функции полезности.

В контрольной политике, которая по постановке не предъявляет участнику никакой рекомендации, индикаторный компонент тождественно равен нулю для всех докладов слота. В этом случае функция полезности сводится к взвешенной сумме компонента релевантности и компонента социального заражения, и распределение выбора определяется содержательной близостью докладов к профилю участника и кумулятивным эффектом в когорте. В этом и состоит роль контрольной политики как точки отсчёта.

\paragraph{Социальное заражение как логарифмический сигнал по числу выбравших доклад.}
Компонент социального заражения отражает динамическое усиление выбора в зависимости от уже сделанных выборов в когорте: формально это монотонно возрастающая функция от числа участников, уже выбравших данный доклад к моменту принятия решения текущим участником. На этапе проектирования рассматривался ряд форм этой зависимости --- линейная пропорциональная по числу выбравших, насыщающая функция с экспоненциальным затуханием, нормированная доля внутри слота, нормировка относительно среднего по слоту, логарифмическая шкала, форма на основе распространения информационных каскадов и несколько производных от перечисленных.

В принятой реализации использована форма с логарифмической нормировкой числа выбравших к логарифму общего размера аудитории. Эта форма сочетает два полезных свойства: монотонный рост влияния с числом ранее сделанных выборов в когорте и насыщение при больших значениях, что предотвращает доминирование первого популярного доклада за счёт эффекта <<победитель забирает всё>>. Она не вводит в модель новых параметров, кроме веса самого компонента, и встречается в литературе по агентским симуляторам пользователей рекомендательных систем.

Значение компонента социального заражения вычисляется по числу выбравших данный доклад, а не по доле занятых мест в зале. Это архитектурно существенно. На текущих программах конференций, где в одном временном слоте каждый зал содержит ровно один доклад, две величины численно совпадают. Семантически они различны: число выбравших доклад --- это популярность доклада, доля занятых мест в зале --- ограничение вместимости. Первая величина даёт понятийное разделение каналов <<социальный сигнал>> и <<ограничение вместимости>> и поддерживает принятое решение о выносе ограничения вместимости в политику. При программах с другой структурой, где в одном зале в одном слоте может проходить несколько докладов, разделение сохраняется без переписывания модели.

В отличие от компонента следования рекомендации, компонент социального заражения вычисляется детерминированно по фиксированному состоянию счётчика выборов и не использует генератора случайных чисел. Это свойство существенно для воспроизводимости сравнения политик: при нулевом весе компонента социального заражения функция полезности сводится к двум остальным компонентам, и вся стохастическая логика прогона совпадает с реализацией без социального заражения с точностью до пословного протокола шагов. Тем самым требование совместимости с базовой моделью при выключенном социальном заражении выполняется по построению, не за счёт отдельной логики.

\paragraph{Распределение выбора и альтернатива отказа от посещения.}
Распределение вероятностей выбора над докладами слота получается применением функции softmax к полученным значениям полезности с делением на параметр температуры. Дополнительно в распределение явно введена альтернатива отказа от выбора с фиксированной вероятностью; остальная вероятностная масса распределяется между докладами слота. Альтернатива отказа моделирует наблюдаемую на конференциях долю участников, не посещающих текущий слот по причинам, не связанным с программой; явное введение <<no-choice alternative>> соответствует распространённой практике моделирования индивидуального выбора в литературе по дискретному выбору.

\paragraph{Простота модели и её сценарная роль.}
Принятая модель отклика проста по форме и интерпретируема по компонентам: три аддитивных канала, симплексное соотношение весов, общая логит-форма распределения. Эта простота продиктована содержательным ограничением задачи: данные о фактической посещаемости отсутствуют, и калибровка модели на эмпирических наблюдениях за индивидуальным поведением невозможна. Введение более тонкой структуры --- гетерогенности параметров по латентным классам, вложенной логит-формы с группированием близких альтернатив, моделей с динамической памятью на уровне отдельного участника --- означает добавление в модель свободы, которая не может быть подтверждена данными и превращается в источник трудно интерпретируемых артефактов. Допущение о параметрической форме модели выбора, на котором держится эта простота, зафиксировано в главе~\ref{ch:method}.

Принятая логит-форма удобна по нескольким причинам. Граничная верификация по весу компонента следования рекомендации выполняется по построению: при нулевом весе все политики становятся неразличимы по функции полезности. Режим общих случайных чисел сохраняется при попарном сравнении политик внутри точки гиперкуба. Состав компонентов функции полезности можно расширять, не переписывая её общей структуры; так в реализации был добавлен компонент социального заражения после фиксации основной двухкомпонентной модели на этапе ранней проверки.

Сценарная роль модели отклика отделена от роли прогноза. Модель не претендует на воспроизведение фактического поведения аудитории конкретной конференции; она описывает фиксированную параметризуемую структуру отклика, относительно которой выполняется численное сравнение политик и вариантов программы. Чувствительность результатов сравнения к допущениям модели исследуется в эксперименте отдельно через варьирование параметрических осей, описывающих веса компонентов, размер аудитории и источник оценки популярности докладов. Эта проверка устойчивости --- самостоятельная задача, выполняемая в главе~\ref{ch:results}.


\section{Реализация управляющих воздействий: политики рекомендаций и оператор локальных модификаций программы}
\label{sec:3-4}

У организатора программы конференции имеется два качественно различных канала влияния на ожидаемое распределение посещений. Первый канал --- алгоритмический: формирование рекомендаций отдельным участникам в каждом временном слоте через выбранную политику. Второй канал --- структурный: локальная модификация привязки докладов к временным слотам с сохранением общего состава программы. В программном комплексе оба канала реализованы как взаимозаменяемые компоненты, подключаемые к симулятору через единый контракт; настоящий подраздел описывает реализацию каждого из них и обоснование выбора конкретных дизайнов.

\paragraph{Семейство политик рекомендаций.}
В сравнении участвует семейство из четырёх политик, отвечающее минимальной достаточности для содержательного сравнительного анализа.

Контрольная политика не предъявляет участнику никакой рекомендации. В функции полезности участника компонент следования рекомендации тождественно равен нулю на всех докладах слота, и распределение выбора определяется только содержательной близостью и компонентом социального заражения. Эта политика выполняет роль точки отсчёта: при любых значениях параметрических осей её результаты задают опорную картину распределения посещений, формирующегося без алгоритмического вмешательства. Без такой точки отсчёта сравнение остальных политик теряло бы интерпретируемость, поскольку любой эффект политики был бы неотличим от эффекта самой по себе аддитивной функции полезности.

Политика по релевантности формирует упорядоченный список рекомендаций по убыванию содержательной близости описания доклада к профилю участника. Сигналы загрузки залов и популяции она не использует. Эта политика моделирует наиболее распространённый в практических рекомендательных системах малого масштаба класс подходов, в том числе тот, который реализован в существующем сервисе рекомендаций партнёра. Она служит ориентиром для оценки эффекта учёта системных ограничений в более сложных политиках.

Политика с учётом загрузки залов сочетает ранжирование по релевантности с явным учётом текущей заполненности залов. К скорингу кандидатов добавляется компонент, монотонно убывающий по доле занятых мест в зале, в котором проходит доклад; залы с заполненностью выше заданного жёсткого порога исключаются из выдачи. При недостаточном числе допустимых кандидатов жёсткое отсечение снимается, чтобы выдача сохраняла ненулевую длину. Эта политика реализует простое содержательное правило, выдвигаемое организатором мероприятия: при прочих равных направлять участников в менее загруженные залы и не направлять в почти переполненные. В ней сосредоточен основной канал учёта вместимости, вынесенный из функции полезности участника по соображениям, изложенным в подразделе~\ref{sec:3-3}.

Политика на основе языковой модели формирует выдачу через вызов внешней языковой модели, которой передаётся текстовый профиль участника и текстовые описания кандидатов слота с просьбой вернуть упорядоченный список идентификаторов. Эта политика реализует четвёртый качественно различный класс --- семантическое ранжирование без явно заданной функции близости. От политики по релевантности она отличается тем, что строит ранжирование через семантический разбор текстов и не опирается на фиксированную косинусную близость векторных представлений; от политики с учётом загрузки залов --- отсутствием явного канала ограничений вместимости. Включение её в семейство проверяет, даёт ли семантическое ранжирование заметный выигрыш относительно простой косинусной близости при равенстве остальных условий. В реализации использован вариант без передачи в запрос состояния загрузки залов: это сохраняет вынос ограничения вместимости на уровень соответствующей политики семейства и обеспечивает сопоставимость с политикой по релевантности по составу учитываемых сигналов. Результаты вызовов кэшируются по идентификаторам участника и слота для воспроизводимости и контроля повторных запусков.

Минимальная достаточность четырёх политик обоснована тем, что каждая представляет качественно различный класс подхода: отсутствие алгоритмического вмешательства, классическое ранжирование по релевантности, учёт системных ограничений вместимости, семантическое ранжирование на основе языковой модели. Обучаемые политики на основе графовых нейронных сетей или обучения с подкреплением исключены из семейства из-за отсутствия обучающего материала, что обсуждалось в главах~\ref{ch:domain} и~\ref{ch:method}. Алгоритмы переранжирования с управлением разнообразием выдачи рассматривались отдельно: задача управления разнообразием не сводится к учёту вместимости и не входит в центральный вопрос работы --- сравнение политик по показателям риска перегрузки.

Все политики реализованы как взаимозаменяемые компоненты с единым контрактом подключения к симулятору. Это даёт прозрачное сравнение политик в одной экспериментальной обстановке и возможность расширения семейства без изменения симулятора и компонентов оценки. Политика на основе языковой модели подключается через тот же контракт, что и три остальные, с единственным отличием --- стоимостью и латентностью вызовов; этот аспект учитывается в дизайне эксперимента, не в логике симулятора.

\paragraph{Оператор локальных модификаций программы.}
Оператор локальной модификации программы реализует структурный канал управления. Он принимает на вход исходную программу и возвращает конечное множество допустимых вариантов программы, полученных перестановками пар докладов между разными временными слотами одного дня. Модификация состоит в том, что у двух выбранных докладов меняется временной слот; привязка к залу при этом сохраняется, состав докладов программы остаётся неизменным. Полученные варианты подаются в симулятор так же, как исходная программа, и образуют отдельную ось параметрической конфигурации.

При выборе формы оператора рассматривался ряд альтернатив. Полная задача оптимизации программы средствами целочисленного программирования, разработанная в литературе по conference scheduling --- в работах Вангервена и соавторов, Резайнии и соавторов, Пилявского и соавторов, --- была отвергнута по двум причинам. Во-первых, такие постановки опираются на допущение известных или восстановленных опросом преференций участников; в условиях настоящей работы соответствующий канал данных недоступен. Во-вторых, и это главное, задача работы --- не нахождение оптимального расписания, а сценарный анализ устойчивости программы к её локальным изменениям. Применение полной оптимизации в этой роли было бы методологически некорректным: сравнение исходной программы с глобально оптимизированной не отражает интересующего вопроса об эффекте небольших, реалистичных для организатора перестановок.

Альтернативные локальные операторы --- обмен блоков по нескольким докладам, перестановки между залами внутри одного слота, эвристический локальный поиск по риску перегрузки --- также рассматривались. Обмен блоков увеличивает мощность изменения и приближает оператор к полной оптимизации, что снижает интерпретируемость результата как <<небольшой реалистичной перестановки>>. Перестановки между залами внутри одного слота не меняют распределение посещений по слотам и интересны как отдельный самостоятельный канал, выходящий за рамки настоящей работы. Эвристический локальный поиск по риску перегрузки превращает оператор в оптимизатор и противоречит роли структурной оси сценарного эксперимента.

Принятая форма оператора --- попарная перестановка двух докладов между разными временными слотами одного дня. Это минимальная локальная модификация, не выходящая за пределы одного дня программы, понятная организатору как реалистичный приём управления программой и дающая равномерную сопоставимость вариантов. Бюджет оператора задаётся числом возвращаемых модификаций; при превышении бюджета модификации выбираются случайной подвыборкой с фиксированным случайным зерном на уровне точки гиперкуба, что даёт воспроизводимость и одинаковость множества вариантов для всех политик внутри одной точки.

Ограничение на конфликты спикеров реализовано как жёсткая проверка: модификация, в результате которой один спикер оказывается одновременно в двух разных временных слотах одного дня, исключается из множества возвращаемых вариантов. Проверка применяется ко всем кандидатам и предшествует случайной подвыборке. На программах, в исходных данных которых поле спикеров отсутствует, проверка становится тривиальной: результаты в этом случае трактуются как отсутствие данных о возможных конфликтах, а не как отсутствие конфликтов в качестве свойства программы. Это уточнение существенно для интерпретации результатов на наборах данных без поля спикеров и фиксируется как одно из условий применения оператора.

Исходная программа включается в множество вариантов как отдельная контрольная точка с фиксированным индексом. При сравнении эффекта структурных модификаций исходная программа всегда присутствует на той же оси параметрической конфигурации, что и её модификации, и попарное сравнение модификаций с исходной программой выполняется без отдельной служебной логики.

Алгоритмический и структурный каналы управления реализованы как ортогональные оси параметрической конфигурации: при каждом запуске эксперимента симулятор получает сначала вариант программы, выбранный по соответствующей оси, затем политику рекомендаций, выбранную по другой оси, и обрабатывает их совместно. Это позволяет в главе~\ref{ch:results} раздельно оценить эффект алгоритмической логики и эффект структуры программы на показатели риска и устойчивости.


\section{Реализация LLM-агентского симулятора}
\label{sec:3-5}

В составе комплекса реализован второй симулятор отклика участника на основе большой языковой модели. Он не заменяет параметрический симулятор и не служит самостоятельным результатом работы; его роль --- взаимная проверка выводов параметрического симулятора о сравнительной устойчивости политик на репрезентативном подмножестве конфигураций. Параметрический и агентский симуляторы образуют не два независимых в строгом смысле источника отклика, а два источника отклика с принципиально разной формой модели выбора: в параметрическом --- фиксированная функция полезности с явными весами компонентов, в агентском --- свободная форма решения через цепь рассуждений языковой модели по тексту профиля и контекста. Совпадение или расхождение выводов двух источников при общем входе и общем контракте политики служит диагностическим сигналом о чувствительности результата к форме модели выбора. Настоящий подраздел описывает реализацию агента, состав передаваемого ему контекста и обоснование принятых ограничений.

\paragraph{Состав агента.}
Агент реализован в минимальной архитектуре <<профиль --- память --- действие>>. Профиль агента --- это текстовое описание участника, согласованное с тематикой конференции; в нём фиксируются интересы, уровень квалификации, предпочтения по форматам выступлений и стилевые ожидания. Память агента --- это упорядоченный список докладов, посещённых им в текущем эпизоде до момента принятия очередного решения; при формировании очередного решения список передаётся агенту в текстовой форме с указанием идентификаторов и заголовков посещённых докладов. Действие агента --- выбор одного из доступных в текущем временном слоте докладов или явный отказ от посещения слота; решение возвращается в структурированной форме с кратким текстовым обоснованием.

В текущем эпизоде агенты обрабатываются последовательно по слотам и последовательно внутри одного слота. Каждый следующий агент в слоте наблюдает обновлённое состояние, отражающее решения тех агентов, что были обработаны раньше; так формируется кумулятивное распределение посещений, согласованное с параметрическим симулятором. После завершения слота агент фиксирует посещение в собственной памяти, и эта запись становится доступной ему при принятии решений в последующих слотах того же эпизода.

Эта архитектура --- стандартная минимальная сборка для агентских симуляторов пользователей рекомендательных систем, использованная в работе Парка и соавторов о генеративных агентах и в средах Agent4Rec, OASIS и AgentSociety, обзор которых приведён в главе~\ref{ch:domain}. Расширения минимальной архитектуры --- пятифакторная модель личности, граф знакомств между агентами, рефлексивная переработка опыта, межслотовое общение участников через общий пул впечатлений --- рассматривались на этапе проектирования и оставлены за пределами объёма работы. Каждое такое расширение добавляет поведенческие элементы, которые сами требуют отдельной проверки на согласованность; роль LLM-симулятора в настоящей работе не в более точном моделировании аудитории, а в независимом источнике отклика для перекрёстной проверки выводов параметрического симулятора на общем подмножестве конфигураций.

\paragraph{Состав передаваемого контекста и паритет с параметрическим симулятором.}
Существенная часть проектирования LLM-агента --- определение того, что именно передаётся в запрос модели в качестве контекста принятия решения. Принципом, по которому делалось это определение, является паритет каналов с параметрическим симулятором: по каждому каналу влияния на выбор участника решается, передавать соответствующую информацию агенту или нет, и решение это принимается симметрично с реализацией соответствующего канала в функции полезности параметрического симулятора.

Текстовый профиль и история сегодняшних посещений передаются всегда: они соответствуют каналу содержательной релевантности и динамике индивидуального предпочтения по ходу эпизода. Упорядоченный список рекомендаций политики передаётся всегда: он соответствует каналу следования рекомендации и позволяет сравнивать политики в LLM-симуляторе на тех же основаниях, что и в параметрическом. Состояние загрузки залов агенту не передаётся: по принятому в подразделе~\ref{sec:3-3} разделению каналов ограничение вместимости трактуется как ответственность политики, и передача его в промпт агента нарушила бы это разделение, дав агенту прямой доступ к информации, которая в параметрическом симуляторе ему недоступна.

Социальный сигнал --- число участников, уже выбравших каждый из докладов слота к моменту принятия решения текущим агентом --- передаётся в промпт условно. При выключенном канале социального заражения, когда соответствующий вес в функции полезности параметрического симулятора равен нулю, социальный сигнал не передаётся и в LLM-симулятор. При ненулевом весе сигнал передаётся в виде кратких числовых сводок, и в системный промпт добавляется указание учитывать его как умеренный или заметный социальный фактор в зависимости от значения веса. Это решение выровнено с реализацией социального заражения в параметрическом симуляторе: оба источника отклика работают с одной и той же семантикой социального сигнала. Конкретные числовые значения сигнала формируются эндогенно внутри каждого симулятора по ходу эпизода и могут различаться между симуляторами; различие трактуется как следствие различной формы модели выбора, а не как несоответствие реализаций.

\paragraph{Бюджет и ограничения масштаба.}
LLM-симулятор требует существенно больших вычислительных ресурсов и материальных затрат, чем параметрический. Каждое решение каждого агента требует отдельного вызова внешней языковой модели; стоимость и время выполнения растут пропорционально произведению размера аудитории, числа слотов и числа сравниваемых политик. Применить LLM-симулятор на полном плане параметрического эксперимента нецелесообразно ни по времени, ни по бюджету.

В принятой схеме LLM-симулятор работает на репрезентативном подмножестве конфигураций, отобранных по принципу равномерного покрытия параметрического пространства, и на сокращённой синтетической аудитории. Внутри одной точки этого подмножества все четыре политики прогоняются на одной и той же синтетической аудитории --- реализация принципа общих случайных чисел, аналогичная параметрическому симулятору. Контроль стоимости реализован как явное ограничение сверху на суммарные затраты прогона: при достижении заданного предела прогон останавливается, частичные результаты сохраняются, и факт превышения предела фиксируется в метаданных эксперимента. Все вызовы языковой модели выполняются асинхронно с ограничением числа одновременных запросов, для каждого вызова фиксируется фактическая стоимость и время выполнения.

\paragraph{Принципиальная роль второго симулятора.}
Назначение LLM-симулятора в работе --- не получить более достоверную модель аудитории, а проверить выводы параметрического симулятора альтернативным источником отклика. Это совпадает с критической линией литературы по агентским симуляторам, обсуждённой в главе~\ref{ch:domain}: в работе Лароожа и Тёрнберга показано, что индивидуальная точность LLM-агентов невелика, а правдоподобие их поведения не эквивалентно операциональной валидности модели. Поэтому LLM-симулятор не служит источником самостоятельных абсолютных утверждений о поведении аудитории. Все его выводы рассматриваются только в сопоставлении с выводами параметрического симулятора на общем подмножестве конфигураций; совпадение ранжирований политик между двумя симуляторами трактуется как взаимная проверка двух источников отклика с разными формами модели выбора.


\section{Реализация экспериментального дизайна и подготовки данных для валидации}
\label{sec:3-6}

Реализация плана эксперимента отделена от логики симуляторов и оформлена в виде самостоятельного модуля, формирующего набор параметрических конфигураций по заданным осям и обеспечивающего разделение источников случайности между ними. Настоящий подраздел описывает реализацию плана эксперимента и обоснование принятых решений.

\paragraph{Латинский гиперкуб как метод покрытия пространства параметров.}
Пространство параметрических конфигураций задано шестью осями: множитель вместимости залов, источник оценки популярности докладов, вес компонента следования рекомендации, вес компонента социального заражения, размер синтетической аудитории, индекс варианта программы. Третий компонент функции полезности --- содержательная релевантность --- получает дополняющий вес из симплексного ограничения и собственной оси не образует. Параметр стохастичности softmax (температура) и базовая вероятность отказа от посещения слота в плане эксперимента не варьируются: они рассматриваются как фиксированные настройки симулятора, а не как оси основного сценарного эксперимента.

Обоснование выбора латинского гиперкуба как плана эксперимента приведено в главе~\ref{ch:method}. В реализации план содержит пятьдесят точек по шести осям; внутри каждой точки перебираются все четыре политики семейства, и каждая комбинация <<точка $\times$ политика>> повторяется на трёх независимых случайных зёрнах. Этот объём согласован с временным бюджетом полного параметрического прогона и даёт достаточную статистическую базу для попарного сравнения политик внутри точки в режиме общих случайных чисел. Полный факторный перебор по тому же пространству при разумной дискретизации даёт количество комбинаций, превышающее временной бюджет в десятки раз; глобальные методы оценки чувствительности избыточны для задачи попарного сравнения политик внутри точки конфигурации.

\paragraph{Симплексное ограничение весов и сбалансированность дискретных осей.}
Веса трёх компонентов функции полезности участника связаны симплексным ограничением: их сумма равна единице. При генерации точек гиперкуба не все случайные комбинации значений соответствующих координат удовлетворяют этому ограничению, и точки, нарушающие его, исключаются. Реализован циклический алгоритм генерации блоков точек с последующей фильтрацией: каждый блок латинского гиперкуба порождается детерминированно от собственного под-зерна, образованного от главного зерна эксперимента; точки, нарушающие симплекс, отбрасываются; цикл продолжается до набора требуемого числа допустимых точек. Такой порядок воспроизводим и не зависит от порядка генерации блоков.

После накопления требуемого числа точек выполняется проверка сбалансированности дискретных осей --- индекса варианта программы, размера аудитории и источника оценки популярности докладов. Поскольку рандомизация в латинском гиперкубе сама по себе не гарантирует сбалансированности дискретных осей при небольшом числе точек, при недостаточном представлении какого-либо уровня дискретной оси выполняется замена случайной избыточно представленной точки на принудительно представленную точку соответствующего уровня. Замены фиксируются в метаданных плана и не нарушают свойство равномерного покрытия гиперкуба, поскольку выполняются на оси, дискретной по построению.

\paragraph{Подмножество для агентского симулятора.}
Агентский симулятор работает на репрезентативном подмножестве из двенадцати точек, отобранных из полных пятидесяти. Принципом отбора служит равномерное покрытие пространства: подмножество должно охватывать как крайние, так и центральные значения каждой из осей конфигурации. Этому требованию отвечает алгоритм максимизации минимального попарного расстояния --- последовательный отбор точек, при котором каждая следующая точка максимально удалена от уже отобранных в нормированном пространстве осей. Альтернативные методы --- отбор по кластерам, ручной отбор по краям и центру, равномерное проектирование на каждую ось --- рассматривались на этапе проектирования; алгоритм максимизации минимального расстояния выбран как метод, формализуемо реализующий требование равномерного покрытия без участия эксперта.

Дополнительно к алгоритму отбора применяется принудительное включение точки с исходным вариантом программы: это сохраняет в подмножестве контрольную точку структурного канала управления и поддерживает интерпретацию исходной программы как опорной для сравнения её локальных модификаций.

\paragraph{Общие случайные числа.}
Принцип общих случайных чисел применяется на уровне точки гиперкуба. Внутри одной точки случайное зерно отбора синтетической аудитории и случайное зерно оператора локальных модификаций программы фиксированы; они одинаковы для всех политик и для всех реплик внутри этой точки. Все политики и все реплики работают с одной и той же синтетической аудиторией и с одной и той же реализацией варианта программы. Допустимый источник случайности при репликации --- собственное зерно симулятора --- изолирован от первых двух и варьируется между репликами. Внутри самого симулятора разделены поток случайности окончательного выбора участника и поток случайности стохастической части политики; это делает граничное поведение модели строго воспроизводимым.

Такое разделение существенно для содержательной интерпретации сравнения политик. Без общих случайных чисел шум, вносимый разной выборкой аудитории при сравнении двух политик, мог бы маскировать или искусственно усиливать различие между ними. С общими случайными числами это различие отражает только различие в логике политик при одинаковом внешнем условии; именно это и нужно для попарного сравнения внутри точки конфигурации.

\paragraph{Ограничение запуска политики на основе языковой модели.}
Стоимость запуска политики на основе языковой модели существенно превосходит стоимость запуска остальных политик семейства. Полный охват ею всех пятидесяти точек параметрического плана эксперимента не укладывается в практический временной и денежный бюджет работы. В принятой схеме политика на основе языковой модели запускается в параметрическом симуляторе только на тех же двенадцати точках, что и подмножество для агентского симулятора. Это даёт перекрёстную сопоставимость двух симуляторов на общих двенадцати точках и сохраняет возможность сравнения трёх остальных политик семейства на полных пятидесяти точках. Решение фиксируется в метаданных эксперимента и в разделе ограничений работы.

\paragraph{Формат результатов.}
Результаты единичных итераций сохраняются в табличном формате с явным указанием идентификатора точки гиперкуба, значений всех параметрических осей конфигурации, идентификатора политики, номера реплики и значений всех показателей групп, описанных в подразделе~2.2.3. Формат напрямую совместим с модулем постобработки, который вычисляет попарные сравнения политик, оценки устойчивости при варьировании одной из осей и эффект структурных модификаций программы. Постобработка не зависит от внутренней логики симуляторов: она работает как функция, отображающая таблицу результатов в таблицы попарных сравнений и характеристик чувствительности. Численные результаты этой обработки и их интерпретация представлены в главе~\ref{ch:results}.


\section{Использованные данные и верификация работоспособности}
\label{sec:3-7}

Программный комплекс работает на наборе данных, состоящем из программы основной конференции, контролируемых синтетических программ для проверки граничных свойств реализации и пред-сгенерированной популяции синтетических профилей участников. Настоящий подраздел описывает назначение каждого типа данных и реализованные процедуры проверки работоспособности комплекса.

\paragraph{Программа основной конференции.}
В качестве основной конференции для численного эксперимента использована программа Mobius 2025 Autumn --- мероприятия сообщества JUG, посвящённого мобильной разработке. Программа охватывает два дня и содержит сорок докладов, распределённых по шестнадцати временным слотам: в двенадцати параллельных слотах используются три зала, в четырёх пленарных --- один. Номинальная вместимость каждого зала составляет сто посадочных мест. Выбор этой программы обусловлен сочетанием трёх факторов: наличие у партнёра реальной программы с фактическими данными о залах и докладах, размер программы, достаточный для содержательного сценарного анализа без чрезмерной вычислительной стоимости, и согласованность тематики программы с тематикой подготовленной популяции синтетических профилей участников.

Метаданные программы --- текстовые описания докладов, тематические признаки и состав спикеров --- используются на разных этапах работы комплекса. Текстовые описания служат основой для формирования векторных представлений докладов, используемых в политиках рекомендаций и в функции полезности модели поведения участника. Тематические признаки и сопутствующие структурные характеристики программы используются для вычисления оценки популярности докладов в одном из вариантов источника популярности. Состав спикеров используется оператором локальных модификаций программы для проверки ограничения на конфликты спикеров.

\paragraph{Синтетическая популяция участников.}
Аудитория задаётся синтетическим порождающим процессом. Для основной конференции подготовлена популяция из ста индивидуальных профилей участников. Каждый профиль --- это текстовое описание интересов, уровня квалификации, тематических предпочтений и стилевых ожиданий, согласованное с тематикой основной конференции, и нормированное векторное представление этого описания, используемое в функции полезности модели поведения. Синтетические аудитории конкретных размеров формируются отбором подмножеств из этой популяции с фиксированным случайным зерном, образуемым по описанному в подразделе~\ref{sec:3-6} правилу.

Пред-сгенерированная популяция с последующим отбором подмножеств выбрана в противовес генерации новой популяции на каждом запуске. Это сокращает стоимость подготовки данных, даёт воспроизводимость и согласованность аудиторий между прогонами, поддерживает корректное применение принципа общих случайных чисел на уровне точки гиперкуба. Разнородность популяции --- профили разной квалификации и разных тематических подобластей --- фиксируется на этапе подготовки и не зависит от случайного зерна отбора подмножества.

\paragraph{Контролируемые синтетические программы.}
Помимо программы основной конференции в составе комплекса подготовлены три контролируемые синтетические программы --- toy-микроконференции, --- назначение которых отлично от роли основной конференции. Эти программы используются для верификации граничных свойств реализации и не участвуют в основном сценарном эксперименте.

Первая toy-микроконференция содержит один временной слот с двумя параллельными залами и двумя докладами и предназначена для проверки базовых свойств функции полезности и распределения выбора при контролируемых параметрах: при нулевом весе компонента следования рекомендации все политики дают близкие распределения посещений, при ненулевом весе различие политик выражено и монотонно растёт по этому весу. Вторая toy-микроконференция содержит два временных слота с такой же структурой каждого слота и предназначена для проверки кумулятивного эффекта между слотами в LLM-агентском симуляторе: на одной слотовой программе агент с памятью сегодняшних посещений неотличим от агента без памяти, и для содержательной проверки этого канала требуется как минимум двухслотовая программа. Третья контролируемая программа содержит специально внесённый конфликт спикеров --- одного спикера в двух разных временных слотах одного дня --- и предназначена для проверки того, что оператор локальных модификаций программы корректно отвергает варианты, нарушающие соответствующее ограничение.

\paragraph{Граничная верификация модели.}
Граничная верификация модели реализована как блокирующий фильтр перед содержательными выводами в соответствии с четырьмя свойствами, сформулированными в главе~\ref{ch:method}: исчезающий риск перегрузки при избыточной номинальной вместимости, монотонный рост риска при её уменьшении, неразличимость политик при нулевом весе компонента следования рекомендации и максимальное различие политик при единичном весе того же компонента.

Проверка реализована как набор автоматизированных тестов, каждый из которых выполняется на минимально достаточном размере конфигурации с усреднением по нескольким случайным зёрнам. Тесты предусмотрены для запуска при проверке изменений ядра комплекса; при провале хотя бы одного из них переход к содержательным результатам блокируется до устранения причины. Такая процедура соответствует канону верификации имитационных моделей [Sargent, 2013; Robinson, 2014], упомянутому в главе~\ref{ch:domain}. На текущей редакции комплекса все четыре граничных свойства выполняются --- соответствующая часть автоматизированных тестов проходит без отказов; численные характеристики граничной верификации представлены в главе~\ref{ch:results}.

\paragraph{Юнит-тестовое покрытие модулей.}
Помимо граничной верификации в составе комплекса реализованы юнит-тесты на отдельные модули: вычисление компонентов функции полезности, вычисление показателей загрузки и риска, корректность работы политик, корректность работы оператора локальных модификаций программы и проверки ограничения на конфликты спикеров, корректность работы генератора плана эксперимента, корректность правил образования случайных зёрен. Юнит-тесты не пересекаются с граничной верификацией: первые проверяют корректность отдельных компонентов в изоляции, вторая --- корректность модели в целом на уровне эмерджентных свойств. Два уровня тестирования вместе позволяют проверить, что отдельные модули комплекса работают корректно, а собранная из них модель ведёт себя ожидаемо в граничных режимах.


\section{Технологический стек и использование систем искусственного интеллекта}
\label{sec:3-8}

\paragraph{Технологический стек.}
Технологическая сложность реализованного комплекса определяется не выбором конкретных библиотек, а инженерными решениями, описанными в подразделах~\ref{sec:3-3},~\ref{sec:3-5} и~\ref{sec:3-6}: разделением каналов влияния политики и состояния системы; двусимуляторной архитектурой моделирования отклика с паритетом каналов между двумя симуляторами; реализацией принципа общих случайных чисел через явное разделение источников случайности на трёх уровнях (отбор аудитории, оператор локальных модификаций программы, симулятор) и двух потоков случайности внутри симулятора (выбор участника, стохастическая часть политики); оператором локальных модификаций программы с жёсткой проверкой ограничений; асинхронной обработкой вызовов внешних языковых моделей с контролем стоимости и кэшированием. Конкретный технологический стек собран из типовых библиотек экосистемы Python; нестандартная сложность вынесена в перечисленные архитектурные решения.

Программный комплекс реализован на языке Python. Численные вычисления выполняются средствами библиотек NumPy и SciPy; SciPy используется в том числе для генерации точек латинского гиперкуба. Обучение вспомогательной модели предпочтений, использованной в одном из вариантов оценки релевантности, выполнено средствами библиотеки scikit-learn методом градиентного бустинга на решающих деревьях. Векторные представления описаний докладов и текстовых профилей участников получены однократным запросом модели многоязычных эмбеддингов \texttt{intfloat/multilingual-e5-small}; результаты сохранены в бинарных файлах с привязкой к идентификаторам.

Доступ к большим языковым моделям, использованным в политике на основе языковой модели и в LLM-агентском симуляторе, реализован через единый совместимый с OpenAI протокол с маршрутизацией через сервис OpenRouter. В политике на основе языковой модели использована модель \texttt{openai/gpt-4o-mini}; в LLM-агентском симуляторе для основного прогона на репрезентативном подмножестве конфигураций использована экономичная модель того же семейства; конкретное имя модели зафиксировано в метаданных каждого запуска. Все вызовы выполняются асинхронно с ограничением числа одновременных запросов и обработкой ошибок внешнего сервиса; для каждого вызова фиксируется фактическая стоимость и время выполнения. Реализован дисковый кэш ответов политики на основе языковой модели --- повторный запуск той же конфигурации не оплачивается дважды; ответы LLM-агентов в текущей реализации не кэшируются.

Управление виртуальным окружением и зависимостями выполнено с использованием утилиты uv. Хранение версий кода и результатов экспериментов организовано через систему контроля версий git с ведением единой основной ветки. Метаданные каждого запуска эксперимента --- точные значения параметров, используемые случайные зёрна, время выполнения, фактическая стоимость в случае использования внешних языковых моделей --- сохраняются в составе результатов; этого достаточно для воспроизведения параметрической части эксперимента и трассировки условий LLM-запусков.

\paragraph{Использование систем искусственного интеллекта при подготовке работы.}
В соответствии с требованиями к выпускной квалификационной работе, действующими в Университете ИТМО, в составе работы фиксируется применение систем искусственного интеллекта при её подготовке.

Системы искусственного интеллекта применялись в работе как инструмент моделирования аудитории внутри программного комплекса, как инструмент разработки и как инструмент подготовки текстов и анализа источников. Первое применение --- внутренний компонент разрабатываемой системы и описано в подразделах~\ref{sec:3-4} и~\ref{sec:3-5}: большая языковая модель работает как ранжирующий компонент в одной из политик семейства и как источник отклика в агентском симуляторе. Это применение --- часть предмета работы и подвергается тем же процедурам верификации, что и остальные модули комплекса.

Применение систем искусственного интеллекта в качестве инструмента разработки и инструмента подготовки текстов выполнялось вспомогательно. Языковые модели использовались для генерации фрагментов кода по техническим заданиям, формулируемым автором, для предложения вариантов оформления отдельных текстовых разделов, для пояснения деталей внешних библиотек и для поиска и реферирования научных источников при подготовке обзора предметной области в главе~\ref{ch:domain}. Полученные фрагменты во всех случаях проходили содержательный контроль автора: для кода --- функциональное тестирование и интеграцию в комплекс с прохождением реализованных юнит-тестов и граничной верификации; для текстов --- редакторскую проверку на соответствие постановке работы и на стилистическое единство с остальной частью текста; для источников --- независимую проверку существования и содержания каждой цитируемой работы по первичным источникам. Принятие архитектурных решений, постановка задачи, выбор методов, интерпретация полученных результатов и итоговая ответственность за работу принадлежат автору.


\section{Выводы по главе}
\label{sec:3-9}

В главе представлена реализация программного комплекса для сценарного анализа программы конференции в условиях неопределённости спроса и отсутствия данных о фактической посещаемости. Состав комплекса отвечает постановке главы~\ref{ch:method} и собран из взаимозаменяемых компонентов: внутреннее представление программы конференции, синтетическая аудитория, модель поведения участника, семейство политик рекомендаций, оператор локальных модификаций программы, два варианта симулятора отклика, генератор плана эксперимента, модуль показателей и средства формирования сравнительного отчёта.

Структуру комплекса определяют три инженерных решения, выполненных в ходе работы. Первое --- разделение каналов влияния политики и состояния системы на функцию полезности участника, при котором ограничение вместимости вынесено из модели поведения на уровень одной из политик семейства; благодаря этому требования граничной верификации выполняются по построению. Второе --- двусимуляторная архитектура моделирования отклика: параметрический симулятор работает на полном плане эксперимента, а агентский симулятор на основе языковой модели --- на репрезентативном подмножестве конфигураций для перекрёстной проверки выводов. Третье --- план эксперимента на основе латинского гиперкуба с общими случайными числами внутри точки конфигурации; это даёт равномерное покрытие пространства параметров и низкошумное попарное сравнение политик.

Сравнительные результаты численного эксперимента на программе Mobius 2025 Autumn, попарные оценки устойчивости политик при варьировании параметрических осей, эффект локальных модификаций программы относительно исходной и согласованность ранжирования политик между параметрическим и LLM-агентским симуляторами на общем подмножестве конфигураций представлены в главе~\ref{ch:results}.

\chapter{ВАЛИДАЦИЯ, РЕЗУЛЬТАТЫ ЭКСПЕРИМЕНТОВ И АНАЛИЗ}
\label{ch:results}
\sloppy


\section{Цель и схема экспериментальной валидации}
\label{sec:4-1}

Постановка задачи и метод сценарной оценки риска перегрузки залов сформулированы в главе~\ref{ch:method}; реализующий их программный комплекс описан в главе~\ref{ch:implementation}. Настоящая глава предъявляет результаты применения комплекса к программе конференции Mobius 2025 Autumn и содержит экспериментальную валидацию принятых архитектурных решений.

Валидация в работе понимается в смысле, принятом в литературе по имитационному моделированию: она не утверждает, что построенная модель воспроизводит фактическое поведение аудитории конкретной конференции, поскольку для такого утверждения необходимы данные о фактической посещаемости, в постановке работы недоступные. Валидация состоит в проверке трёх свойств комплекса. Первое: модель сохраняет обязательные граничные свойства --- при свободной вместимости риск перегрузки исчезает, при сжатой растёт; при нулевом весе компонента следования рекомендации политики неразличимы по показателям загрузки, при единичном весе различимы максимально. Второе: относительные сравнения политик и вариантов программы остаются устойчивыми при варьировании параметрических осей. Третье: проверяется степень согласованности ранжирования политик в параметрическом и LLM-агентском симуляторах на общем подмножестве конфигураций.

Все три свойства проверяются на едином плане эксперимента: пятидесяти точках латинского гиперкуба с тремя независимыми случайными зёрнами на каждую комбинацию <<точка $\times$ политика>>, репрезентативном подмножестве из двенадцати точек для агентского симулятора и фиксированных значениях параметров стохастичности и базовой вероятности отказа от посещения. Полный объём прогонов составляет 486 единичных итераций параметрического симулятора и 48 единичных итераций LLM-агентского симулятора.

Структура главы соответствует логической последовательности проверок. Подраздел~\ref{sec:4-2} предъявляет результаты граничной верификации модели как блокирующего фильтра перед содержательными выводами. Подраздел~\ref{sec:4-3} описывает результаты полного параметрического прогона --- общее распределение показателей по политикам и подмножество конфигураций, в которых риск перегрузки физически возникает. Подраздел~\ref{sec:4-4} раскрывает попарные сравнения политик и устойчивость ранжирования. Подраздел~\ref{sec:4-5} рассматривает эффект локальных модификаций программы. Подраздел~\ref{sec:4-6} исследует чувствительность результата к параметрическим осям. Подраздел~\ref{sec:4-7} посвящён согласованности двух симуляторов. Подраздел~\ref{sec:4-8} фиксирует ограничения эксперимента и угрозы валидности. Подраздел~\ref{sec:4-9} подводит итог.

Все численные значения, приведённые в главе, получены из зафиксированных артефактов полного экспериментального прогона на программе Mobius 2025 Autumn. Сводные результаты представлены в таблицах~\ref{tab:ec-verification}--\ref{tab:cross-validation}; состав артефактов эксперимента приведён в приложении~\ref{app:artefacts}. Конкретные технические детали реализации не повторяются: они описаны в главе~\ref{ch:implementation}.


\section{Верификация граничных свойств модели}
\label{sec:4-2}

Граничная верификация модели --- блокирующий фильтр: до её прохождения сравнения политик и вариантов программы не интерпретируются содержательно. Четыре граничных свойства из подраздела~2.4.4 главы~\ref{ch:method} проверяются автоматизированными тестами на минимально достаточных конфигурациях с усреднением по нескольким случайным зёрнам.

Свойство первое. При сильно избыточной номинальной вместимости --- множитель не менее 3.0 относительно номинальной --- среднее относительное превышение вместимости тождественно равно нулю независимо от выбранной политики. Это базовая внутренняя согласованность: при заведомо достаточной вместимости перегрузка не может возникнуть ни при какой стратегии распределения участников.

Свойство второе. При уменьшении номинальной вместимости показатели риска перегрузки не убывают. Проверяется на сетке множителей вместимости с усреднением по пяти независимым случайным зёрнам. Тест считается пройденным, если среднее значение показателя на каждом следующем (более сжатом) уровне не ниже значения на предыдущем (более свободном) уровне.

Свойство третье. При нулевом весе компонента следования рекомендации различие политик по показателям загрузки исчезает. Реализованная функция полезности при нулевом весе не зависит от рекомендации, и распределение выбора участника определяется только содержательной близостью и компонентом социального заражения. Этот инвариант контролируется тестом на пословное совпадение протоколов прогона при разных политиках при одинаковых случайных зёрнах: при нулевом весе все политики порождают один и тот же протокол с точностью до записи сигналов о рекомендациях.

Свойство четвёртое. При единичном весе компонента следования рекомендации различие политик по показателям загрузки максимально. Проверяется на стрессовой конфигурации вместимости: размах значений показателя между политиками в этом режиме существенно превосходит коэффициент вариации того же показателя в режиме нулевого веса.

В составе комплекса реализованы десять автоматизированных тестов: четыре теста соответствуют перечисленным свойствам, остальные шесть расширяют их на конфигурации с ненулевым весом компонента социального заражения. Сводный результат граничной верификации представлен в таблице~\ref{tab:ec-verification}.

\begin{small}
\begin{longtable}{|>{\raggedright\arraybackslash}p{3.2cm}|>{\raggedright\arraybackslash}p{4.2cm}|>{\raggedright\arraybackslash}p{5.5cm}|>{\centering\arraybackslash}p{1.6cm}|}
\caption{Результаты граничной верификации модели}\label{tab:ec-verification}\\
\hline
\textbf{Проверка} & \textbf{Проверяемое свойство} & \textbf{Ожидаемое поведение} & \textbf{Результат} \\
\hline
\endfirsthead
\caption*{Продолжение таблицы~\ref{tab:ec-verification}}\\
\hline
\textbf{Проверка} & \textbf{Проверяемое свойство} & \textbf{Ожидаемое поведение} & \textbf{Результат} \\
\hline
\endhead
\hline
\endfoot
\hline
\endlastfoot
EC1 & Свободная вместимость & При множителе вместимости $\geq 3.0$ среднее относительное превышение вместимости равно нулю для всех политик & Пройдено \\
\hline
EC2 & Монотонность по вместимости & При уменьшении множителя вместимости средний риск не убывает (усреднение по 5 зёрнам) & Пройдено \\
\hline
EC3 & Граница по $w_{rec}$ снизу & При $w_{rec} = 0$ протоколы прогонов разных политик пословно совпадают & Пройдено \\
\hline
EC4 & Граница по $w_{rec}$ сверху & При $w_{rec} = 1$ размах показателя между политиками существенно превосходит шум репликации & Пройдено \\
\hline
EC1 (расширение) & EC1 при $w_{gossip} = 0.3$ & Превышение вместимости остаётся нулевым при включённом социальном заражении и свободной вместимости & Пройдено \\
\hline
EC2 (расширение) & EC2 при $w_{gossip} = 0.3$ & Монотонность по вместимости сохраняется при включённом социальном заражении & Пройдено \\
\hline
EC3 (расширение) & EC3 при $w_{gossip} > 0$ & Совпадение протоколов прогонов сохраняется при включённом социальном заражении & Пройдено \\
\hline
EC4 (расширение) & EC4 при ненулевом $w_{gossip}$ & На границе симплекса при $w_{rec} = 0.7$ и $w_{gossip} = 0.3$ различие политик по показателям загрузки сохраняется & Пройдено \\
\hline
Сохранение эффекта политики с учётом загрузки & Сохранение эффекта политики с учётом загрузки залов при $w_{gossip} > 0$ & При $w_{gossip} = 0.3$ и стрессовой вместимости политика с учётом загрузки залов не уступает политике по релевантности по риску перегрузки & Пройдено \\
\hline
Монотонность концентрации & Усиление концентрации выбора по $w_{gossip}$ & Коэффициент Джини загрузки залов не убывает по сетке $w_{gossip} \in \{0,\ 0.3,\ 0.7\}$ & Пройдено \\
\end{longtable}
\end{small}

Граничная верификация пройдена полностью: 10 из 10 автоматизированных тестов завершились без отказов. Поэтому дальнейшие сравнения политик выполняются на реализации, удовлетворяющей обязательным граничным свойствам модели.

Прохождение граничной верификации означает, что реализованная функция полезности и схема разделения источников случайности удерживают разделение каналов влияния, заявленное в постановке: без этого условия попарные сравнения политик не имели бы методического смысла, поскольку политики не были бы инструментально отделены от состояния системы. Кроме того, модель не порождает контринтуитивных артефактов на границах параметрических осей. Содержательные сравнения политик и вариантов программы в остальных подразделах главы опираются именно на эту проверенную реализацию.

Граничная верификация не утверждает, что модель воспроизводит реальное поведение аудитории; такое утверждение при отсутствии данных о фактической посещаемости невозможно по построению. Она утверждает только, что обязательные свойства, без которых сценарный анализ был бы некорректен, в реализации не нарушены.


\section{Результаты параметрического сценарного эксперимента}
\label{sec:4-3}

Полный параметрический прогон выполнен на программе Mobius 2025 Autumn по плану, описанному в подразделе~\ref{sec:3-6}. Внутри каждой из пятидесяти точек гиперкуба перебирались политики с тремя независимыми случайными зёрнами; политика на основе языковой модели в параметрическом режиме запускалась только на двенадцати точках репрезентативного подмножества по соображениям стоимости. Суммарный объём --- 486 единичных итераций. Главный аналитический интерес представляет подмножество тех точек гиперкуба, в которых риск перегрузки фактически возникает: на остальных точках любое сравнение политик вырождается, поскольку при свободной вместимости все политики дают нулевой риск независимо от логики ранжирования.

Распределение центральных показателей по политикам в полном прогоне на пятидесяти точках представлено в таблице~\ref{tab:central-metrics}.

\begin{table}[h!]
\caption{Центральные показатели параметрического прогона по политикам ($n = 50$ точек гиперкуба, 3 реплики на точку)}
\label{tab:central-metrics}
\centering
\small
\begin{tabular}{|l|r|r|r|r|}
\hline
\textbf{Политика} & \textbf{Среднее превышение} & \textbf{Медиана} & \textbf{P75} & \textbf{Средняя релевантность} \\
\hline
Контрольная               & 0.0395 & 0.0000 & 0.0000 & 0.7308 \\
\hline
По релевантности          & 0.0449 & 0.0000 & 0.0066 & 0.7321 \\
\hline
С учётом загрузки залов   & 0.0336 & 0.0000 & 0.0000 & 0.7314 \\
\hline
\end{tabular}
\end{table}

Сравнение средних в таблице~\ref{tab:central-metrics} даёт первое впечатление о близости трёх политик: различия между ними лежат в пределах сотых долей единицы относительного превышения, а медианы у всех трёх политик равны нулю. Устойчивый сравнительный разрыв между политиками выявляется не средними по выборке, а попарным анализом внутри точки гиперкуба (подраздел~\ref{sec:4-4}), и сосредоточен на risk-positive подмножестве, описанном ниже. Тем не менее уже на уровне сводных характеристик распределения видно: политика с учётом загрузки залов даёт наименьшее среднее относительное превышение вместимости при практически неизменной средней релевантности выбранных докладов. Различие в средней релевантности между тремя политиками не превышает 0.002, то есть переход от чистого ранжирования по релевантности к учёту состояния залов не сопровождается заметной потерей содержательного качества рекомендаций.

Нулевая медиана риска перегрузки в полном прогоне --- это структурное свойство сценарного анализа, а не отсутствие эффекта политик: латинский гиперкуб по построению покрывает всё пространство параметрических конфигураций, в том числе режимы свободной вместимости, в которых риск нулевой по физическим причинам. Из пятидесяти точек гиперкуба ненулевой риск хотя бы у одной из трёх политик контрольного состава возникает в тринадцати. Сводка по этому подмножеству приведена в таблице~\ref{tab:risk-positive}.

\begin{table}[h!]
\caption{Подмножество точек с ненулевым риском перегрузки (risk-positive)}
\label{tab:risk-positive}
\centering
\small
\begin{tabular}{|>{\raggedright\arraybackslash}p{10cm}|r|}
\hline
\textbf{Показатель} & \textbf{Значение} \\
\hline
Всего точек LHS в плане эксперимента & 50 \\
\hline
Точек, на которых хотя бы у одной из политик П1--П3 возникает ненулевой риск & 13 \\
\hline
Доля risk-positive точек & 26~\% \\
\hline
Точек с ненулевым риском у политики по релевантности & 13 / 13 \\
\hline
Точек с ненулевым риском у контрольной политики & 11 / 13 \\
\hline
Точек с ненулевым риском у политики с учётом загрузки залов & 10 / 13 \\
\hline
Точек с ненулевым риском у политики на основе языковой модели (на 12 maximin) & 2 / 12 \\
\hline
\end{tabular}
\end{table}

Содержательный анализ риска должен в первую очередь смотреть на risk-positive подмножество: на остальных точках все политики дают нулевой риск по физическим причинам, и сравнение политик там вырождается. Различие частот ненулевого риска между контрольной политикой и политикой по релевантности отражает структурный эффект ранжирования по содержательной близости: при отсутствии учёта состояния залов рекомендация по релевантности концентрирует поток участников на одних и тех же темах, что в части конфигураций приводит к перегрузке одного из залов.

Среди risk-positive точек выделяется подмножество критически непреодолимых конфигураций --- точек, в которых размер аудитории превышает суммарную вместимость залов в параллельном слоте при текущем множителе вместимости. Таких конфигураций три: при множителях вместимости 0.534, 0.963 и 0.774 в сочетании с размером аудитории сто. Физическое переполнение в них неустранимо никакой стратегией распределения участников. Точки оставлены в анализе для прозрачности: они задают границу применимости сценарного сравнения и позволяют организатору отделить случаи, в которых рекомендация бессильна, от случаев, в которых политика даёт выигрыш.

Полная таблица единичных итераций сохранена в формате long-table с явным указанием значений всех осей конфигурации, идентификатора политики, номера реплики и значений всех показателей групп, описанных в подразделе~2.2.3. Эта таблица служит исходным материалом для всех последующих процедур постобработки, описанных в подразделах~\ref{sec:4-4}--\ref{sec:4-7}.


\section{Попарное сравнение политик и устойчивость выводов}
\label{sec:4-4}

Сравнение политик в работе ведётся попарно внутри одной точки гиперкуба, а не средним показателем по всему пространству конфигураций. Этот выбор содержательно обусловлен тем, что среднее по пространству конфигураций скрывает сценарную чувствительность результата, а в задаче с глубокой неопределённостью именно чувствительность представляет интерес.

Для каждой пары политик в каждой точке гиперкуба сравниваются три значения центральных показателей --- по три независимых случайных зёрна на политику --- и фиксируется, какая политика лидирует. Сравнение выполняется в двух режимах: строгое неравенство и неравенство с допуском. Допуск выбран равным 0.005 для показателей риска и неравенства загрузки и 0.001 для показателя релевантности; эти значения выбраны на этапе постобработки и зафиксированы в таблице $\varepsilon$. Введение допуска позволяет различить случаи фактического превосходства одной политики над другой и случаи практически равных результатов в пределах статистического шума репликации. Доли точек, в которых политика выигрывает строго, выигрывает с допуском, не отличается с допуском, проигрывает строго и проигрывает с допуском, в сумме дают единицу с точностью до округления.

Сводный результат попарного сравнения трёх контрольных политик в полном прогоне для метрики среднего относительного превышения вместимости при допуске $\varepsilon = 0.005$ приведён в таблице~\ref{tab:pairwise-overload}.

\begin{table}[h!]
\caption{Попарное сравнение политик по риску перегрузки в полном прогоне (50 точек, $\varepsilon = 0.005$)}
\label{tab:pairwise-overload}
\centering
\small
\begin{tabularx}{\linewidth}{|>{\raggedright\arraybackslash}p{3.6cm}|>{\centering\arraybackslash}p{1.7cm}|>{\raggedright\arraybackslash}p{3.0cm}|>{\centering\arraybackslash}p{1.5cm}|>{\centering\arraybackslash}p{1.5cm}|>{\raggedright\arraybackslash}X|}
\hline
\textbf{Пара политик} & \textbf{Равны с $\varepsilon$} & \textbf{Победитель за пределами допуска} & \textbf{Доля побед} & \textbf{Доля побед второй} & \textbf{Вывод} \\
\hline
Контрольная vs по релевантности & 86~\% & Контрольная & 14~\% & 0~\% & Контрольная превосходит политику по релевантности в 14~\% точек; обратная победа за пределами допуска не зафиксирована \\
\hline
Контрольная vs с учётом загрузки залов & 86~\% & С учётом загрузки залов & 14~\% & 0~\% & Политика с учётом загрузки залов превосходит контрольную политику в 14~\% точек; обратная победа за пределами допуска не зафиксирована \\
\hline
По релевантности vs с учётом загрузки залов & 78~\% & С учётом загрузки залов & 22~\% & 0~\% & Политика с учётом загрузки залов превосходит политику по релевантности в 22~\% точек; обратная победа за пределами допуска не зафиксирована \\
\hline
\end{tabularx}
\end{table}

Наиболее устойчивый результат попарного сравнения: политика по релевантности не выигрывает у политики с учётом загрузки залов ни в одной из 50 точек гиперкуба ни строго (\texttt{win\_strict = 0}), ни за пределами допуска $\varepsilon = 0.005$ (\texttt{win\_eps = 0}). Политика с учётом загрузки залов, напротив, выигрывает у политики по релевантности в 26~\% точек строго и в 22~\% точек за пределами допуска. На risk-positive подмножестве из тринадцати точек политика с учётом загрузки залов не уступает максимуму из контрольной политики и политики по релевантности на всех 13 / 13 точках в пределах допуска и строго снижает риск перегрузки относительно этого максимума в 11 / 13 точек ($\approx$~85~\%).

Шесть наибольших снижений риска при переходе от политики по релевантности к политике с учётом загрузки залов представлены в таблице~\ref{tab:risk-reductions}.

\begin{table}[h!]
\caption{Наибольшие снижения риска при переходе от политики по релевантности к политике с учётом загрузки залов (risk-positive подмножество)}
\label{tab:risk-reductions}
\centering
\small
\begin{tabular}{|r|r|r|r|r|r|}
\hline
\textbf{Точка LHS} & \textbf{Множитель} & \textbf{Аудитория} & \textbf{Риск (П2)} & \textbf{Риск (П3)} & \textbf{Снижение} \\
\hline
49 & 0.774 & 100 & 0.5449 & 0.3558 & 0.1891 \\
\hline
26 & 0.629 &  60 & 0.1706 & 0.0040 & 0.1667 \\
\hline
35 & 0.963 & 100 & 0.0606 & 0.0025 & 0.0581 \\
\hline
 3 & 0.534 & 100 & 1.2222 & 1.1759 & 0.0463 \\
\hline
36 & 0.696 &  60 & 0.1424 & 0.1111 & 0.0312 \\
\hline
18 & 1.040 &  60 & 0.0214 & 0.0000 & 0.0214 \\
\hline
\end{tabular}
\end{table}

Самое крупное снижение зарегистрировано на точке 49: одна из трёх критически непреодолимых конфигураций (audience > effective per-slot capacity), на которой политика с учётом загрузки залов снижает риск с 0.5449 до 0.3558 при том, что физическое переполнение полностью неустранимо. Точка 26 показывает крупное абсолютное снижение риска: с 0.1706 до 0.0040, то есть почти до нулевого значения.

Устойчивость результата отдельной политики на отдельной точке проверяется средствами анализа разброса показателей между репликами одной комбинации <<точка $\times$ политика>>. На полном прогоне трёх контрольных политик зафиксированы 24 сочетания <<точка $\times$ политика $\times$ показатель>> из 600, в которых отношение стандартного отклонения по репликам к абсолютному значению среднего превышает 0.5; на двенадцати точках с политикой на основе языковой модели --- 4 из 36. Все эти сочетания соответствуют режимам редких событий перегрузки, в которых редкая ненулевая реплика порождает большое относительное стандартное отклонение при малом абсолютном значении среднего; центральные выводы попарного сравнения политик в подразделах~\ref{sec:4-3}--\ref{sec:4-4} не зависят от этих конкретных комбинаций. Подробное распределение разброса по репликам сохранено в составе артефактов постобработки (приложение~\ref{app:artefacts}).

В подмножестве из двенадцати точек, на которых запускалась также политика на основе языковой модели, попарные сравнения по показателю риска перегрузки согласуются с полным прогоном. Политика с учётом загрузки залов даёт одинаковые с допуском результаты с контрольной политикой в 83~\% точек и выигрывает за пределами допуска в 17~\%; с политикой по релевантности --- на тех же 83~\% одинаковые с допуском и в 17~\% точек выигрывает за пределами допуска. По показателю среднего относительного превышения вместимости политика с учётом загрузки залов даёт минимальное среднее (0.0093), политика на основе языковой модели --- 0.0109, контрольная политика --- 0.0127, политика по релевантности --- 0.0136. Эти средние не сопоставимы со средними по пятидесяти точкам: десять из двенадцати точек подмножества попадают в режимы свободной вместимости, и абсолютные значения смещены вниз. Политика на основе языковой модели присутствует только в этом подмножестве.

Тепловая карта ранжирования политик по четырём показателям на двенадцати точках репрезентативного подмножества приведена на рисунке~\ref{fig:ranking-heatmap}. Цветовая шкала отражает позицию политики в ранжировании: тёмные клетки соответствуют верхним позициям, светлые --- нижним. Видна устойчивость лидерства политики с учётом загрузки залов на большинстве точек по показателям риска при разреженной структуре ранжирований по показателям, на которых параметрический симулятор вырождается.

\begin{figure}[h!]
\centering
\includegraphics[width=0.85\linewidth]{analysis_ranking_heatmap_maximin.png}
\caption{Ранжирование политик по четырём показателям на 12 точках maximin-подмножества}
\label{fig:ranking-heatmap}
\end{figure}


\section{Эффект локальных модификаций программы}
\label{sec:4-5}

Оператор локальных модификаций программы, описанный в подразделе~\ref{sec:3-4}, генерирует допустимые варианты программы перестановками пар докладов между разными временными слотами одного дня. Индекс варианта программы фигурирует как одна из шести осей параметрической конфигурации; в полном прогоне на пятидесяти точках индекс варианта принимает шесть уровней, соответствующих исходной программе и пяти её локальным модификациям. Поскольку выборка по плану эксперимента строится латинским гиперкубом, разные уровни индекса варианта программы попадают в разные точки гиперкуба, не выровненные между собой по остальным пяти осям. Это техническое свойство выборки имеет принципиальное методическое следствие: попарные сравнения <<исходная программа против модификации $k$>> при фиксированных значениях остальных осей в текущей выборке не выполнены, и любые агрегаты по уровню индекса варианта программы на полной выборке носят диагностический характер, а не утверждают причинной разности между уровнями.

Диагностические агрегаты среднего относительного превышения вместимости по уровням индекса варианта программы в полном прогоне для политики по релевантности приведены в таблице~\ref{tab:program-effect}.

\begin{table}[h!]
\caption{Диагностические агрегаты по индексу варианта программы (политика по релевантности)}
\label{tab:program-effect}
\centering
\small
\begin{tabularx}{\linewidth}{|r|r|r|>{\raggedright\arraybackslash}X|}
\hline
\textbf{Индекс варианта} & \textbf{Число точек LHS} & \textbf{Среднее превышение} & \textbf{Интерпретация} \\
\hline
0 (исходная программа) &  7 & 0.0099 & Контрольный уровень \\
\hline
1 &  9 & 0.0211 & Заметное изменение в большую сторону \\
\hline
2 &  9 & 0.0633 & Заметное изменение в большую сторону \\
\hline
3 &  8 & 0.0000 & Все точки данного уровня в режиме свободной вместимости \\
\hline
4 &  7 & 0.1762 & Наибольшее агрегатное превышение в выборке \\
\hline
5 & 10 & 0.0183 & Близко к контрольному уровню \\
\hline
\end{tabularx}
\end{table}

Таблица не доказывает причинный эффект конкретной перестановки: уровни индекса программы в плане эксперимента не сопоставлены по остальным осям LHS. Знаковый тест разности между уровнями 0 и 4 на cross-product парах даёт p-значение около 0.84; гипотеза об отсутствии разности на текущей выборке не отвергается. Размах агрегата между уровнями 3 и 4 (от 0.0000 до 0.1762) отражает прежде всего различие остальных осей вокруг разных уровней индекса программы, а не различие самих модификаций.

Результат раздела состоит в том, что оператор $\Phi$ интегрирован в сценарный эксперимент как полноценная параметрическая ось наравне с вместимостью, размером аудитории и весами компонентов модели поведения; сводки по уровням этой оси (таблица~\ref{tab:program-effect}) дают сценарные характеристики распределения показателей риска по вариантам программы. Положение №4 защищаемых результатов закрывается именно как количественные сценарные оценки распределения по оси конфигурации, а не как причинная разность между конкретными перестановками. Сравнение <<исходная программа против модификации~$k$>> при фиксированных остальных осях относится к локальной оптимизации программы и оставлено за рамками работы; этот шаг зафиксирован в подразделе~\ref{sec:4-8} как направление развития.

Вторая сторона того же результата --- реализованный комплекс позволяет численно сопоставить эффект изменения структуры программы с эффектом изменения политики рекомендаций в одной экспериментальной обстановке. Эта возможность отвечает четвёртому защищаемому положению работы о совместном оценивании вклада алгоритмической логики и структуры программы в показатели риска.


\section{Анализ влияния параметрических осей}
\label{sec:4-6}

Чувствительность результатов к параметрическим осям конфигурации исследуется через сопоставление подмножеств точек гиперкуба, отличающихся значениями одной из осей при варьируемых остальных. Такое сопоставление конфаундировано: варьирующиеся остальные оси вносят собственный вклад в наблюдаемую разность, поэтому каждая сводка трактуется как диагностический сигнал о направлении и порядке эффекта, а не как причинная оценка. Глобальные методы оценки чувствительности с разложением вклада осей рассматривались на этапе проектирования и отвергнуты как избыточные для задачи попарного сравнения политик внутри точки.

Влияние множителя вместимости и веса компонента следования рекомендации в данных согласуется со свойствами граничной верификации (подраздел~\ref{sec:4-2}): множитель меньше единицы концентрирует случаи ненулевого риска (10 из 50 точек), множитель от двух до трёх даёт нулевой риск во всех 19 точках; различие политик растёт по мере увеличения веса компонента следования рекомендации. Эти зависимости не приносят нового содержательного знания и приведены в составе сводок чувствительности (приложение~\ref{app:artefacts}).

Распределение среднего относительного превышения вместимости по бакетам веса социального заражения представлено в таблице~\ref{tab:gossip-buckets}.

\begin{table}[h!]
\caption{Средний риск перегрузки по бакетам веса социального заражения $w_{gossip}$}
\label{tab:gossip-buckets}
\centering
\small
\begin{tabular}{|c|r|r|r|r|}
\hline
\textbf{Бакет $w_{gossip}$} & \textbf{Число точек} & \textbf{Контрольная} & \textbf{По релевантности} & \textbf{С учётом загрузки залов} \\
\hline
$[0,\ 0.25)$  & 23 & 0.0061 & 0.0123 & 0.0003 \\
\hline
$[0.25,\ 0.5)$ & 19 & 0.0891 & 0.0957 & 0.0822 \\
\hline
$[0.5,\ 0.7]$ &  8 & 0.0178 & 0.0178 & 0.0139 \\
\hline
\end{tabular}
\end{table}

Графически тот же агрегат показан на рисунке~\ref{fig:gossip-bucket-bar}: высота столбцов отражает средний риск перегрузки по политикам в каждом из трёх бакетов веса социального заражения.

\begin{figure}[h!]
\centering
\includegraphics[width=0.85\linewidth]{analysis_gossip_bucket_bar.png}
\caption{Средний риск перегрузки по бакетам веса социального заражения}
\label{fig:gossip-bucket-bar}
\end{figure}

В данных наблюдается немонотонная связь между весом социального заражения и средним риском: максимум среднего риска приходится на средний бакет $w_{gossip}$. Это диагностический результат, а не причинная оценка: остальные оси LHS в бакетах не фиксированы, и наблюдаемый пик в среднем бакете может объясняться конфаундированностью с осями вместимости и размера аудитории, а не свойствами самого канала социального заражения. Внутри каждого бакета политика с учётом загрузки залов сохраняет минимальное среднее значение риска среди трёх политик контрольного состава.

Влияние размера аудитории и источника оценки популярности докладов в данных умеренное и не порождает новых содержательных закономерностей; соответствующие сводки сохранены в составе сводок чувствительности (приложение~\ref{app:artefacts}).

Связь риска и релевантности на пятидесяти точках исследуется через двумерное распределение пар <<риск перегрузки $\times$ релевантность>>. Сводка по частоте trade-off приведена в таблице~\ref{tab:tradeoff}.

\begin{table}[h!]
\caption{Частота trade-off между снижением риска и потерей релевантности ($\varepsilon_{\text{risk}} = 0.005$, $\varepsilon_{\text{utility}} = 0.001$)}
\label{tab:tradeoff}
\centering
\small
\begin{tabularx}{\linewidth}{|>{\raggedright\arraybackslash}X|r|r|r|}
\hline
\textbf{Подмножество} & \textbf{Всего комбинаций} & \textbf{С trade-off} & \textbf{Доля} \\
\hline
Полный прогон трёх основных политик & 150 & 11 & 7.3~\% \\
\hline
Подмножество с участием LLM-политики &  48 &  3 & 6.3~\% \\
\hline
\end{tabularx}
\end{table}

Двумерное распределение пар <<риск $\times$ релевантность>> по всем точкам и политикам показано на рисунке~\ref{fig:risk-utility-scatter}. Большинство точек сгруппированы в зоне нулевого риска при сопоставимой релевантности; точки trade-off лежат в правой верхней части картины и не образуют доминирующей моды.

\begin{figure}[h!]
\centering
\includegraphics[width=0.85\linewidth]{analysis_risk_utility_scatter.png}
\caption{Распределение пар <<риск перегрузки $\times$ релевантность>> по политикам}
\label{fig:risk-utility-scatter}
\end{figure}

Trade-off редок в обоих подмножествах: на большинстве точек гиперкуба политика, лидирующая по снижению риска, не сопровождает это лидерство заметной потерей содержательного качества рекомендаций. Это согласуется с близкими значениями средней релевантности у трёх политик в таблице~\ref{tab:central-metrics}.


\section{Согласованность параметрического и LLM-агентского симуляторов}
\label{sec:4-7}

Вторая проверка относительного порядка политик выполнена прогоном LLM-агентского симулятора на репрезентативном подмножестве из двенадцати точек гиперкуба, отобранных алгоритмом максимизации минимального попарного расстояния с принудительным включением точки исходной программы. Внутри одной точки все четыре политики прогоняются на одной и той же синтетической аудитории. Объём прогона --- 48 единичных итераций (двенадцать точек $\times$ четыре политики $\times$ одно случайное зерно). Технические характеристики прогона --- фактическая стоимость, время выполнения, число вызовов внешней языковой модели и доля ошибок парсинга --- сохранены в метаданных эксперимента и в основном тексте не используются.

Согласованность ранжирования политик между параметрическим и LLM-агентским симуляторами оценивается по четырём ключевым показателям. Для каждой из двенадцати точек подмножества вычисляются ранжирования политик в обоих симуляторах, и попарное соответствие ранжирований количественно характеризуется коэффициентом ранговой корреляции Спирмена. Минимальная медианная корреляция, при которой согласованность считается принятой, фиксирована на уровне 0.5. Сводный результат представлен в таблице~\ref{tab:cross-validation}.

\begin{table}[h!]
\caption{Согласованность ранжирования политик между параметрическим и LLM-агентским симуляторами на 12 точках подмножества}
\label{tab:cross-validation}
\centering
\small
\begin{tabularx}{\linewidth}{|>{\raggedright\arraybackslash}p{4.0cm}|>{\centering\arraybackslash}p{2.0cm}|>{\centering\arraybackslash}p{1.6cm}|>{\centering\arraybackslash}p{1.6cm}|>{\raggedright\arraybackslash}X|}
\hline
\textbf{Показатель} & \textbf{Невырожденных точек} & \textbf{Медиана $\rho$} & \textbf{Совпадение лидера} & \textbf{Интерпретация} \\
\hline
Средняя релевантность & 12 & 0.80 & 11 / 12 & Высокая согласованность \\
\hline
Доля слотов с переполнением & 2 & 0.74 & 2 / 2 & Диагностическая согласованность из-за малого числа невырожденных точек \\
\hline
Относительное превышение вместимости & 2 & 0.30 & 2 / 2 & Ниже порога $\rho \geq 0.5$, узкая выборка \\
\hline
Дисперсия загрузки залов в слоте & 12 & 0.40 & 11 / 12 & Ниже порога по $\rho$, но высокое совпадение лидера \\
\hline
\end{tabularx}
\end{table}

Графическая сводка коэффициентов ранговой корреляции Спирмена по метрикам перекрёстной проверки приведена на рисунке~\ref{fig:cross-validation-rho}.

\begin{figure}[h!]
\centering
\includegraphics[width=0.85\linewidth]{cross_validation_rho_per_metric.png}
\caption{Коэффициенты ранговой корреляции Спирмена по метрикам перекрёстной проверки}
\label{fig:cross-validation-rho}
\end{figure}

Объединённая медиана коэффициента корреляции по четырём показателям составляет 0.554 при пороге 0.5. Эта агрегатная величина сама по себе диагностически малоинформативна: она усредняет 28 значений с разной долей валидных точек --- по релевантности и по дисперсии загрузки по 12 точек, по показателям переполнения по 2 точки из-за вырождения параметрического ранжирования на 10 из 12 безопасных точек подмножества. Содержательную картину даёт построчная разбивка таблицы~\ref{tab:cross-validation}. Подробные данные по точкам с пометкой <<degenerate>> для каждой точки с вырожденным ранжированием сохранены в составе артефактов перекрёстной проверки (приложение~\ref{app:artefacts}).

Согласованность параметрического и LLM-агентского симуляторов наблюдается по части показателей: уверенно --- по средней релевантности (12 точек) и по верхнему лидерству политик в дисперсии загрузки (11 / 12 точек); по показателям переполнения корреляция оценивается на 2 невырожденных точках из 12 и остаётся диагностической. Интерпретация расхождений по метрикам переполнения ограничена структурой выборки: десять из двенадцати точек параметрически безопасны, поэтому ранжирование по этим метрикам вырождается. На текущем подмножестве нельзя отделить возможное расхождение моделей от эффекта малой доли невырожденных точек. Третьему защищаемому положению работы это отвечает в части реализации двух механизмов отклика и перекрёстной проверки на репрезентативном подмножестве.

Использованные в LLM-прогоне модели зафиксированы в метаданных эксперимента: политика на основе языковой модели работала на \texttt{openai/gpt-4o-mini}, агентский симулятор --- на \texttt{openai/gpt-5.4-nano}. Обе модели выбраны из соображений практической достижимости бюджета и не претендуют на максимально достижимое качество поведенческой модели; это сознательное ограничение, фиксируемое в подразделе~\ref{sec:4-8}.


\section{Ограничения эксперимента и угрозы валидности}
\label{sec:4-8}

Полученные результаты подлежат интерпретации в рамках принятой постановки задачи. Ограничения работы фиксируются как техническое условие задачи, не как защитная оговорка.

\textit{Отсутствие данных о фактической посещаемости.} Аудитория задана синтетическим порождающим процессом. Все численные значения показателей риска, загрузки и релевантности интерпретируются как сравнительные характеристики политик и вариантов программы внутри принятой модели, не как прогноз фактической посещаемости конкретной конференции.

\textit{Одна основная конференция.} Численные результаты полного прогона получены на программе Mobius 2025 Autumn. Перенос относительного порядка политик на программы других мероприятий требует отдельной экспериментальной проверки; универсальность выводов работы для произвольной программы конференции не утверждается.

\textit{Фиксированные параметры стохастичности и базовой вероятности отказа.} Параметр стохастичности softmax-распределения и базовая вероятность отказа от посещения слота в плане эксперимента не варьировались. Чувствительность результатов к этим параметрам в работе не исследована.

\textit{Ограниченный масштаб LLM-прогона.} Политика на основе языковой модели запускалась только на двенадцати точках гиперкуба; LLM-агентский симулятор работал на тех же двенадцати точках с одним случайным зерном. Использована экономичная модель из соображений бюджета --- это сознательное ограничение и направление развития работы.

\textit{Не полностью воспроизводимые ответы LLM-агентов.} Параметрическая часть эксперимента воспроизводима полностью. Ответы LLM-агентов не кэшируются: повторный прогон LLM-симулятора может дать другой набор индивидуальных решений из-за стохастичности языковой модели. Метаданные условий запуска фиксируются.

\textit{Конфаундированность анализа по оси индекса варианта программы.} Точки гиперкуба с разными уровнями индекса варианта программы не сопоставлены между собой по остальным осям. Агрегаты по этой оси (см. таблицу~\ref{tab:program-effect}) --- диагностические, не причинные.

\textit{Узкая выборка для согласованности по метрикам переполнения.} Корреляция Спирмена между двумя симуляторами по метрикам, непосредственно связанным с переполнением, опирается на две невырожденные точки из двенадцати (см. таблицу~\ref{tab:cross-validation}). Вывод о согласованности по этим метрикам остаётся диагностическим до прогонов на конфигурациях с большей долей risk-positive точек.

\textit{Сценарная, а не прогнозная интерпретация.} Полученные значения трактуются как сравнительные внутри принятой модели; утверждения вида <<при политике P фактическое превышение составит N процентов>> в работе не делаются и из результатов не следуют.


\section{Выводы по главе}
\label{sec:4-9}

Содержательный итог экспериментальной валидации программного комплекса формулируется в семи тезисах.

\begin{enumerate}
\item \textbf{Граничная верификация пройдена полностью.} Все четыре свойства, обязательные для корректной модели рассматриваемого класса, на текущей реализации выполняются; 10 из 10 автоматизированных тестов завершились без отказов (таблица~\ref{tab:ec-verification}).

\item \textbf{Политика с учётом загрузки залов снижает средний риск при сохранении релевантности.} В полном прогоне она даёт меньшее среднее относительное превышение вместимости (0.0336), чем контрольная политика (0.0395) и политика по релевантности (0.0449), при практически неизменной средней релевантности --- различие в пределах 0.002 (таблица~\ref{tab:central-metrics}).

\item \textbf{В попарном сравнении политика по релевантности не выигрывает у политики с учётом загрузки залов ни в одной из 50 точек гиперкуба} ни строго (\texttt{win\_strict = 0}), ни за пределами допуска $\varepsilon = 0.005$ (\texttt{win\_eps = 0}). Обратное направление выигрыша зафиксировано в 26~\% точек строго и в 22~\% точек за пределами допуска (таблица~\ref{tab:pairwise-overload}).

\item \textbf{На risk-positive подмножестве политика с учётом загрузки залов устойчиво лидирует.} В 13 точках с ненулевым риском хотя бы у одной из политик она не уступает максимуму из двух остальных в 13 / 13 точек в пределах допуска и строго снижает риск в 11 / 13 точек (таблицы~\ref{tab:risk-positive} и~\ref{tab:risk-reductions}).

\item \textbf{Trade-off между снижением риска и потерей релевантности редок.} На полном прогоне он встречается в 7.3~\% комбинаций, на подмножестве с участием LLM-политики --- в 6.3~\% (таблица~\ref{tab:tradeoff}).

\item \textbf{Согласованность параметрического и LLM-агентского симуляторов наблюдается по части показателей.} По средней релевантности --- медиана $\rho = 0.80$ на 12 точках и совпадение верхнего лидера на 11 / 12 точках. По дисперсии загрузки залов --- совпадение верхнего лидера на 11 / 12 точках при медиане $\rho = 0.40$ ниже формального порога. По метрикам переполнения параметрическое ранжирование вырождается на 10 из 12 точек, и Spearman $\rho$ оценивается на 2 невырожденных точках; вывод по этим метрикам остаётся диагностическим (таблица~\ref{tab:cross-validation}). Эти результаты частично поддерживают третье защищаемое положение: реализованы два механизма отклика и проведена перекрёстная проверка, но по метрикам переполнения согласованность остаётся диагностической.

\item \textbf{Результаты --- сценарное сравнительное утверждение, не прогноз.} Полученные значения относятся к сравнительному анализу политик и вариантов программы в рамках принятой модели на программе Mobius 2025 Autumn. Они не утверждают воспроизведения фактического поведения аудитории, не выделяют единственной оптимальной политики рекомендаций и не претендуют на универсальность для произвольной программы конференции.
\end{enumerate}

Содержательное значение полученного отчёта для организатора программы конференции состоит в следующем: использование сравнительного отчёта позволяет заранее идентифицировать сценарии параметрической неопределённости, в которых программа уязвима к перегрузке залов, и принять решение о выборе политики рекомендаций или о структурной модификации программы до проведения мероприятия --- на основании сравнительной устойчивости показателей риска внутри принятой модели, а не на основании прогноза фактической посещаемости.

Полное содержание сформированного сравнительного отчёта --- карта ожидаемой загрузки залов, перечень горячих точек программы, попарные сравнения политик, диагностика по уровню индекса варианта программы, анализ чувствительности по параметрическим осям и блок согласованности двух симуляторов --- сохранено в составе результатов и предъявляется как итоговый артефакт системы поддержки принятия решений для организатора программы конференции.

\chapter*{ЗАКЛЮЧЕНИЕ}
\addcontentsline{toc}{chapter}{ЗАКЛЮЧЕНИЕ}

Заключение работы оформляется по итогам экспериментальной части и приводится в финальной версии выпускной квалификационной работы.


% Список использованных источников.
% На текущей стадии ссылки в тексте оформлены как [Author, Year];
% для отображения списка литературы используется \nocite{*}.
\addcontentsline{toc}{chapter}{СПИСОК ИСПОЛЬЗОВАННЫХ ИСТОЧНИКОВ}
\nocite{*}
\printbibliography[title={СПИСОК ИСПОЛЬЗОВАННЫХ ИСТОЧНИКОВ}]

% Приложения. Шаблонное приложение А (appA.tex) намеренно не подключается:
% единственное содержательное приложение оформлено в appB.tex и при включении
% через \appendix получает букву <<А>> в нумерации заголовка.
\appendix
\chapter{СОСТАВ ЭКСПЕРИМЕНТАЛЬНЫХ АРТЕФАКТОВ}
\label{app:artefacts}
\sloppy

В таблице~\ref{tab:artefacts} приведён состав закоммиченных артефактов, на которых получены численные результаты главы~\ref{ch:results}. Артефакты сгруппированы по назначению и хранятся в каталоге \texttt{experiments/results/} репозитория проекта.

\begin{small}
\begin{longtable}{|>{\raggedright\arraybackslash}p{3.0cm}|>{\raggedright\arraybackslash}p{5.8cm}|>{\raggedright\arraybackslash}p{5.7cm}|}
\caption{Состав экспериментальных артефактов}\label{tab:artefacts}\\
\hline
\textbf{Группа} & \textbf{Назначение} & \textbf{Артефакты} \\
\hline
\endfirsthead
\caption*{Продолжение таблицы~\ref{tab:artefacts}}\\
\hline
\textbf{Группа} & \textbf{Назначение} & \textbf{Артефакты} \\
\hline
\endhead
\hline
\endfoot
\hline
\endlastfoot
Таблица единичных прогонов & Long-format таблица всех 486 единичных итераций параметрического симулятора с явным указанием значений всех параметрических осей конфигурации, идентификатора политики, номера реплики и значений показателей групп, описанных в подразделе~2.2.3 & \texttt{lhs\_parametric\_mobius\_2025\_autumn\_2026-05-08.\{json,csv,md\}} \\
\hline
Файлы постобработки параметрического прогона & Сводки попарных сравнений, агрегатов по уровням параметрических осей, диагностики risk-positive подмножества, чувствительности по осям, разброса между репликами, диагностики политики на основе языковой модели и распределения <<риск~$\times$~релевантность>> & \texttt{analysis\_pairwise.json}, \texttt{analysis\_capacity\_audit.json}, \texttt{analysis\_program\_effect.json}, \texttt{analysis\_gossip\_effect.json}, \texttt{analysis\_sensitivity.json}, \texttt{analysis\_stability.json}, \texttt{analysis\_llm\_ranker\_diagnostic.json}, \texttt{analysis\_risk\_utility.json} \\
\hline
Артефакты LLM-агентского симулятора & Long-format таблица 48 единичных итераций LLM-агентского симулятора на 12 точках репрезентативного подмножества & \texttt{llm\_agents\_lhs\_subset\_12pts.\{json,csv,md\}} \\
\hline
Артефакты перекрёстной проверки & Сводка коэффициентов ранговой корреляции Спирмена и совпадений верхнего лидера политик между параметрическим и LLM-агентским симуляторами на 12 общих точках, с пометками вырождения по точкам & \texttt{analysis\_cross\_validation.\{json,md\}} \\
\hline
Сводные отчёты & Компактные технические отчёты по результатам параметрического прогона (стадия~U) и совместного прогона с перекрёстной проверкой (стадия~W); содержат сводные таблицы и приёмочные чек-листы & \texttt{report\_mobius\_2025\_autumn\_parametric.md}, \texttt{report\_mobius\_2025\_autumn\_full.md} \\
\hline
Визуализации & Графические сводки результатов: диаграмма рассеяния <<риск~$\times$~релевантность>>, столбчатая диаграмма среднего риска по бакетам веса социального заражения, тепловая карта ранжирования политик на 12 maximin-точках, диаграмма коэффициентов ранговой корреляции по метрикам перекрёстной проверки & \texttt{plots/analysis\_risk\_utility\_scatter.png}, \texttt{plots/analysis\_gossip\_bucket\_bar.png}, \texttt{plots/analysis\_ranking\_heatmap\_maximin.png}, \texttt{plots/cross\_validation\_rho\_per\_metric.png} \\
\hline
Метаданные LLM-прогонов & Параметры запуска агентского симулятора, конкретные имена использованных языковых моделей (\texttt{openai/gpt-4o-mini} для политики на основе языковой модели, \texttt{openai/gpt-5.4-nano} для агентского симулятора), фактические значения суммарной стоимости и времени выполнения, число вызовов и доля ошибок парсинга структурированного ответа & Поля \texttt{params} и блок acceptance в \texttt{llm\_agents\_lhs\_subset\_12pts.json} \\
\hline
Тесты граничной верификации & Автоматизированные тесты, реализующие граничные свойства EC1--EC4 и их расширения на конфигурации с ненулевым весом социального заражения; каждый тест выполняется на минимально достаточной конфигурации с усреднением по нескольким случайным зёрнам & \texttt{experiments/tests/test\_extreme\_conditions.py} \\
\end{longtable}
\end{small}

Сохранённые артефакты позволяют воспроизвести любую единичную итерацию параметрического прогона по фиксированному значению параметрических осей и случайных зёрен. Воспроизводимость LLM-агентских прогонов ограничена стохастичностью внешней языковой модели и фиксируется на уровне условий запуска.


\end{document}
